% =====================================================================
% 第五讲 (Exposé V) 基本群: 一般理论
% (Le groupe fondamental: généralités)
% A. Grothendieck
% 原书页码: pp. 87-115
% =====================================================================
\chapter{基本群：一般理论}

\section*{0. 引言}

本讨论班是 1960 年讨论班的续篇. 我们用诸如 \ref{cor:I.9.7} 这样的记号引用后者，它表示：代数几何讨论班，第 I 讲，第 9.7 号. 1961 年各讲的编号紧随 1960 年的编号. 我们用诸如 (EGA I 8.7.3) 这样的记号引用 Dieudonné-Grothendieck 的《代数几何原理》(Éléments de Géométrie Algébrique).

本讲概括 (并稍加补充) 1960 年最后几个未成文的讲.

与 1961 年一样，一般限于局部诺特预概形，尽管这个限制常常无关紧要. 我们将在讲 VI 中承认忠实平坦下降理论，它概括在 Bourbaki 讨论班 N$^{\circ}$ 190 中. 若有必要，我们将在以后的讲中给出更详细的阐述\footnote{参见讲 VI 和讲 VIII.}，一旦读者有机会相信这一技术对基本群理论的有用性.

\section{带有限算子群的预概形。商预概形}

设 $X$ 是预概形，$G$ 是通过自同构作用在 $X$ 上的有限群 (为确定起见作用在右边). 若 $X$ 是环为 $A$ 的仿射概形，则 $G$ 通过左自同构作用在 $A$ 上.

对每个预概形 $Z$，$G$ 左作用在集合 $\mathrm{Hom}(X, Z)$ 上，于是可以考虑集合
\[
\mathrm{Hom}(X, Z)^{G}
\]
(被 $G$ 不变的态射). 它函子性地依赖于 $Z$，可以问这个函子是否``可表示'' (representable)，i.e. 同构于函子 $Z \mapsto \mathrm{Hom}(Y, Z)$.

这意味着可以找到预概形 $Y$ 和 $G$-不变态射
\[
p \colon X \longrightarrow Y,
\]
使得对每个 $Z$，相应的映射 $g \mapsto gp$
\[
\operatorname{Hom}(Y, Z) \longrightarrow \operatorname{Hom}(X, Z)^{G}
\]
双射. 此时称 $(Y, p)$ 是 $X$ 除以 $G$ 的商预概形 (quotient prescheme) (它在唯一同构意义下确定).

\begin{proposition}\label{prop:V.1.1}
设 $A$ 是有限群 $G$ 在其上左作用的环，$B = A^{G}$ 是 $A$ 的不变元子环，$X = \operatorname{Spec}(A)$，$Y = \operatorname{Spec}(B)$，$p \colon X \to Y$ 是 (显然 $G$-不变的) 典范态射. 则
\begin{enumerate}[(i)]
\item $A$ 在 $B$ 上整，i.e. $p$ 是整态射 (integral morphism)；
\item $p$ 满射，其纤维是 $G$ 的轨道，$Y$ 的拓扑是 $X$ 的拓扑的商；
\item 设 $x \in X$，$y = p(x)$，$G_x$ 是 $x$ 的稳定化子，则 $k(x)$ 是 $k(y)$ 的拟 Galois 代数扩张，且从 $G_x$ 到 $k(x)$ 的 $k(y)$-自同构群 $\operatorname{Gal}\bigl(k(x)/k(y)\bigr)$ 的典范映射满射；
\item $(Y, p)$ 是 $X$ 除以 $G$ 的商预概形.
\end{enumerate}
\end{proposition}

\begin{proof}
陈述 (i)，(ii)，(iii) 在交换代数中众所周知\footnote{参见 N. Bourbaki，Alg. Comm. Chap. 5，§ 1 和 § 2，th. 2.}，只作为备忘，只有关于拓扑的断言例外，它来自下列一般事实 (Cohen-Seidenberg 定理的简单推论)：整态射是闭的 (i.e. 把闭集变为闭集). 立即指出：
\end{proof}

\begin{corollary}\label{cor:V.1.2}
在以上条件下，自然同态
\[
\mathcal{O}_Y \longrightarrow p_{*}(\mathcal{O}_X)^{G}
\]
是同构.
\end{corollary}

\begin{proof}
这由公式
\[
(S^{-1}A)^{G} = S^{-1}(A^{G})
\]
立即得出，该公式对 $B = A^{G}$ 的每个乘法稳定子集 $S$ 成立 (该公式对模也可模块化，更一般地，对平坦的基变换 $A \to A'$ 成立)，把它应用于 $S$ 由 $B$ 的一个元素 $f$ 生成的情形.
\end{proof}

断言 (ii) 和推论 \ref{cor:V.1.2} 容易蕴含 (iv)；更一般地，有：

\begin{proposition}\label{prop:V.1.3}
设 $X$ 是带有限自同构群 $G$ 的预概形，$p \colon X \to Y$ 是 $G$-不变仿射态射，使得 $\mathcal{O}_Y \stackrel{\sim}{\longrightarrow} p_{*}(\mathcal{O}_X)^{G}$. 则 \ref{prop:V.1.1} 的结论 (i)，(ii)，(iii)，(iv) 仍然成立.
\end{proposition}

\begin{proof}
事实上，对 (i)，(ii)，(iii)，可以假设 $Y$ 从而 $X$ 仿射，若 $B$，$A$ 是其环，假设蕴含 $B = A^{G}$，只需应用 \ref{prop:V.1.1}. 对 (iv)，用 (ii) 和 $\mathcal{O}_Y = p_{*}(\mathcal{O}_X)^{G}$.
\end{proof}

\begin{corollary}\label{cor:V.1.4}
在 \ref{prop:V.1.3} 的条件下，对 $Y$ 的每个开集 $U$，$U$ 是 $X|U = p^{-1}(U)$ 除以 $G$ 的商.
\end{corollary}

\begin{proof}
事实上，由 $p$ 诱导的 $p^{-1}(U) \to U$ 满足与 $p$ 相同的假设.
\end{proof}

若现在 $X$ 是 $Z$-预概形且 $G$ 的作用是 $Z$-自同构，则由 (iv)，$Y$ 是 $Z$-预概形. 据此：

\begin{corollary}\label{cor:V.1.5}
为使 $X$ 在 $Z$ 上仿射 (相应地：分离)，必要且充分条件是 $Y$ 如此. 若 $X$ 在 $Z$ 上有限型，则它在 $Y$ 上有限；若还有 $Z$ 局部诺特，则 $Y$ 在 $Z$ 上有限型.
\end{corollary}

\begin{proof}
由于 $X$ 在 $Y$ 上仿射 (更不用说分离)，若 $Y$ 在 $Z$ 上仿射 (相应地：分离)，则 $X$ 亦然. 反过来，假设 $X$ 在 $Z$ 上仿射，证明 $Y$ 仿射：借助 \ref{cor:V.1.4} 可以假设 $Z$ 仿射，归结为证明若 $X$ 仿射则 $Y$ 仿射，这由 \ref{prop:V.1.1} 中 $Y$ 作为 $\operatorname{Spec}(A)^{G}$ 的显式确定得出. 同样，由于 $p \colon X \to Y$ 整故泛闭，且满射，推出若 $X$ 在 $Z$ 上分离则 $Y$ 分离 (一个有待阐明的引理!)；事实上，在交换图
\[
\begin{tikzcd}[row sep=large, column sep=large]
X \times_Z X \arrow[r,"p \times_Z p"] \arrow[d] & Y \times_Z Y \arrow[d] \\
X \arrow[r,"p"] & Y
\end{tikzcd}
\]
中，态射 $X \times_Z X \to Y \times_Z Y$ 闭，故把 $X \times_Z X$ 的 (闭的) 对角线变为 $Y \times_Z Y$ 的闭子集，由于 $p$ 满射它正是后者的对角线. —— 若 $X$ 在 $Z$ 上有限型，则更不用说在 $Y$ 上有限型，故在 $Y$ 上有限 (因为它已在 $Y$ 上整). 再假设 $Z$ 局部诺特，证明 $Y$ 在 $Z$ 上有限型. 借助 \ref{cor:V.1.4} 可以假设 $Z$ 仿射. 由于拓扑空间 $X$ 拟紧且 $p \colon X \to Y$ 满射，$Y$ 也拟紧，故是有限个仿射开集的并，由 \ref{cor:V.1.4} 归结到 $Y$ 仿射从而 $X$ 仿射的情形. 但此时 $X$ 的环 $A$ 是诺特环 $Z$ 的环 $C$ 上的有限型代数，且众所周知 $B = A^{G}$ 也是 $C$ 上的有限型代数 (因为 $A$ 将整，故在 $B$ 的 $C$ 上有限型的子代数 $B'$ 上有限，而 $B'$ 诺特，故 $B$ 也在 $B'$ 上有限，从而在 $C$ 上有限型).
\end{proof}

\begin{corollary}\label{cor:V.1.6}
为使 $X$ 仿射 (相应地：是概形)，必要且充分条件是 $Y$ 如此.
\end{corollary}

\begin{definition}\label{def:V.1.7}
设 $X$ 是有限群 $G$ 在其上右作用的预概形. 称 $G$ 的作用可允许 (admissible)，若存在具有 \ref{prop:V.1.3} 性质的态射 $p \colon X \to Y$ (这意味着 $X/G$ 存在且同构于 $Y$).
\end{definition}

\begin{proposition}\label{prop:V.1.8}
设 $X$ 是有限群 $G$ 在其上右作用的预概形. 为使 $G$ 的作用可允许，必要且充分条件是 $X$ 是 $G$-不变仿射开集的并，或等价地，$G$ 在 $X$ 中的每个轨道都含于一个仿射开集.
\end{proposition}

\begin{proof}
后一个条件显然由前一个蕴含，反过来它也蕴含前一个；因为设 $T$ 是轨道，$U$ 是含它的仿射开集，$U$ 在 $G$ 诸 $g$ 下的变换之交是含 $T$、含于仿射开集 $U$ 的 $G$-稳定开集 $U'$. 由于在 $U$ 中每个有限子集都有仿射开邻域基本组，存在含于 $U'$ 的 $T$ 的仿射开邻域 $V$. 它在 $G$ 诸 $g$ 下的变换是仿射的且含于分离的 $U'$，故其交 $U''$ 是 $G$-不变的仿射开集且含 $T$. —— 据此，\ref{prop:V.1.8} 中所考虑的条件必要，因为取覆盖 $Y$ 的仿射开集 $Y_i$ 的原像 $X_i$ 即可. 它充分，因为此时可以借助 \ref{prop:V.1.1} 构造商 $Y_i = X_i/G$；在每个 $Y_i$ 中，$X_i \cap X_j$ 的像是开集 $Y_{ij}$，由 \ref{cor:V.1.4} 等同于 $X_{ij}/G$，特别由此得到同构 $Y_{ij} \xrightarrow{\sim} Y_{ji}$，使得可以把诸 $Y_i$ 粘合起来构造 $Y$. Serre 宁愿直接构造 $X$ 除以 $G$ 的商拓扑空间 $Y$，在其上放层 $p_{*}(\mathcal{O}_X)^{G}$，并验证 $Y$ 成为预概形，且处于 \ref{prop:V.1.3} 的条件下.
\end{proof}

% 原书编号笔误: 下面这条推论在原文中与上面的定义同为 ``\ref{def:V.1.7}''
\begingroup
\addtocounter{theorem}{-1}%
\renewcommand{\thetheorem}{\thesection.7}%
\begin{corollary}\label{cor:V.1.7sub}
若作用在 $X$ 上的 $G$ 可允许，则 $G$ 的每个子群 $H$ 亦然 (故 $X/H$ 存在).
\end{corollary}

\endgroup

这也可以直接在 \ref{prop:V.1.3} 的情形下验证，注意总可以假设 $X$ 在某个 $Z$ 上仿射，且 $s \in G$ 的元素通过 $Z$-自同构作用 (例如取 $Z = Y$)；事实上：

\begin{corollary}\label{cor:V.1.8}
假设 $X$ 在 $Z$ 上仿射，且 $G$ 的作用是 $Z$-自同构. 则 $G$ 在 $X$ 上的作用可允许. 若 $X$ 由拟凝聚代数层 $\mathcal{A}$ 定义，则 $Y$ 由 $\mathcal{A}$ 的 $G$ 不变元层 $\mathcal{A}^{G}$ 定义.
\end{corollary}

\begin{proposition}\label{prop:V.1.9}
假设 $G$ 在 $X$ 上的作用可允许，且 $X/G = Y$ 是 $Z$ 上的预概形. 考虑基变换态射 $Z' \to Z$，令 $X' = X \times_Z Z'$，$Y' = Y \times_Z Z'$，于是 $G$ 通过结构的转移仍作用在 $X'$ 上，态射 $p' \colon X' \to Y'$ 不变. 若 $Z'$ 在 $Z$ 上平坦，则 $p'$ 仍满足 \ref{prop:V.1.3} 的假设，i.e. $\mathcal{O}_{Y'} \longrightarrow p'_{*}(\mathcal{O}_{X'})^{G}$ 是同构 ($p'$ 无论如何都是仿射的). 故 $G$ 在 $X'$ 上的作用可允许，且
\[
(X/G) \times_Z Z' \approx (X \times_Z Z')/G.
\]
\end{proposition}

\begin{proof}
显然可以假设 $Z = Y$，归结到 $Y$ 和 $Y'$ 仿射的情形. 必须证明：若 $B$ 是 $G$ 在 $A$ 中作用的不变元子环，$B'$ 是在 $B$ 上平坦的 $B$-代数，则 $B'$ 是 $A' = A \otimes_B B'$ 的不变元子代数. 这立即可见，因为正合序列
\[
0 \longrightarrow B \stackrel{i}{\longrightarrow} A \stackrel{j}{\longrightarrow} A^{(G)}
\]
(其中最后一项表示 $A$ 的幂，且 $j(x)$ 是系统 $s \cdot x - x$ ($s \in G$)) 与 $B'$ 张量后仍正合.
\end{proof}

注意平坦性假设对结果的成立是本质的；特别地，若 $Y'$ 是 $Y$ 的闭子预概形 (例如甚至是 $Y$ 的闭点)，$X'$ 是它在 $X$ 中的原像，则 $Y'$ 一般不等同于 $X'/G$. 我们将看到，若 $X$ 在 $Y$ 上平展，则情况仍然如此.

为结束本节，给出一个既方便又平凡的框架. 设 $Y$ 是预概形. 由于在预概形范畴中直和存在，可以对每个集合 $E$ 考虑一族 $(Y_i)_{i \in E}$ (全部等同于 $Y$) 的直和预概形，记为 $Y \times E$. 它由公式
\begin{equation}
\operatorname{Hom}(Y \times E, Z) = \operatorname{Hom}(E, \operatorname{Hom}(Y, Z)) \tag{($*$)}
\end{equation}
刻画，其中第二个 Hom 当然表示从集合 $E$ 到集合 $\mathrm{Hom}(Y, Z)$ 的映射的集合. 有典范态射
\[
Y \times E \longrightarrow Y,
\]
使 $Y \times E$ 成为 $Y$ 上的预概形. 由于纤维积与直和可交换 (在预概形范畴中)，若 $Y$ 是另一个 $Z$ 上的预概形，则对基变换 $Z' \to Z$ 有：
\[
(Y \times E) \times_Z Z' = (Y \times_Z Z') \times E
\]
(该公式在 $Z = Y$ 时尤其有用). 另一方面，由定义平凡地得出
\[
(Y \times E) \times F = Y \times (E \times F) = (Y \times E) \times_Y (Y \times F)
\]
(不过最后一个公式来自上面指出的交换性).

对固定的 $Y$，可以把 $Y \times E$ 看作 $E$ 的函子，取值于 $Y$ 上的预概形，由前面的公式该函子与有限积可交换 (这使每个通常群 $G$ 对应 $Y$ 上的群概形 $Y \times G$，若 $Y$ 有限则它有限，等等). 更一般地，这个函子``左正合''，但我们在此用不到. 它还与直和可交换，并且``右正合''，这从定义公式 (∗) 立即看出. 特别地，若有限群 $G$ 右作用在集合 $E$ 上，则它右作用在 $Y \times E$ 上，且有
\[
(Y \times E)/G = Y \times (E/G),
\]
其中第一个成员的商确实满足 \ref{prop:V.1.9} 的条件 (立即可见).

\section{分解群与惯性群。平展情形}

设 $G$ 是右作用在预概形 $X$ 上的有限群. 若 $x \in X$，称 $x$ 的\emph{分解群} (decomposition group)\index[term]{fenjiequn@分解群} 为 $x$ 的稳定化子 $G_d(x)$. 该群典范地 (左) 作用在剩余域 $k(x)$ 上，其中平凡作用的元素集称为 $x$ 的\emph{惯性群} (inertia group)\index[term]{guanxingqun@惯性群}，记为 $G_i(x)$.

假设 $G$ 在 $X$ 上的作用可允许，$Y$ 是 $Z$ 上的预概形. 固定 $z \in Z$，以及 $k(z)$ 的代数闭扩张 $\Omega$，其超越次数高于诸 $k(x)/k(z)$ 的超越次数，其中 $x$ 是 $z$ 之上的 $X$ 的点. 可以把 $\operatorname{Spec}(\Omega)$ 看作 $Z$-概形，$X$ 在 $\Omega$ 中取值的点对应于 $k(z)$-代数同态 $k(x) \to \Omega$，其中 $x$ 是 $z$ 之上的 $X$ 的点；由于 $\Omega$ 取得足够大，$z$ 之上的每个点 $x$ 都是 $X$ 在 $\Omega$ 中取值的一个点的所在点. 若 $X(\Omega)$ 和 $Y(\Omega)$ 分别表示 $X$ 和 $Y$ 在 $\Omega$ 中取值的点的集合，则有自然映射
\[
X(\Omega) \longrightarrow Y(\Omega),
\]
另一方面，$G$ 作用在 $X(\Omega)$ 上，前面的映射 $G$-不变. 据此，\ref{prop:V.1.1} 的结论 (ii) 和 (iii) 也解释为：前面的映射满射，且把 $Y(\Omega)$ 等同于商 $X(\Omega)/G$. 此外，若 $x$ 是 $a \in X(\Omega)$ 的所在点，则 $a$ 在 $G$ 中的稳定化子正是惯性群 $G_i(x)$. 这一切甚至在不假设 $\Omega$``足够大''时也成立，后一个假设只用于保证可以把 $z$ 之上 $X$ 的每个元素的惯性群刻画为``几何''稳定化子. 例如立即得出：

\begin{proposition}\label{prop:V.2.1}
作基扩张 $Z' \to Z$，由此 $X' = X \times_Z Z'$. 设 $x'$ 是 $X'$ 的点，$x$ 是它在 $X$ 中的像，则有 $G_i(x) = G_i(x')$.
\end{proposition}

\begin{proof}
只需在上述考虑中取 $\Omega$ 为 $k(z')$ 的足够大的扩张 (其中 $z$，$z'$ 是 $x$，$x'$ 在 $Z$，$Z'$ 中的像).
\end{proof}

\begin{proposition}\label{prop:V.2.2}
在 \ref{prop:V.2.2} 的条件下，假设 $Y$ 局部诺特，$X$ 在 $Y$ 上有限. 设 $H$ 是 $G$ 的子群，考虑 $X' = X/H$ (参见 \ref{def:V.1.7})，设 $x \in X$，$x'$ 是它在 $X'$ 中的像，$y$ 是它在 $Y$ 中的像.
\begin{enumerate}[(i)]
\item 若 $H \supset G_d(x)$，则同态 $\mathcal{O}_y \to \mathcal{O}_{x'}$ 在完备化上诱导同构.
\item 若 $H \supset G_i(x)$，则同态 $\mathcal{O}_y \to \mathcal{O}_{x'}$ 平展，i.e. $X'$ 在 $x'$ 处平展于 $Y$ 上.
\end{enumerate}
\end{proposition}

\begin{proof}
设 $Y_1 = \mathrm{Spec}(\widehat{\mathcal{O}_y})$，作基变换 $Y_1 \to Y$，得到在 $Y_1$ 上有限的 $X_1 = X \times_Y Y_1$，$G$ 在其上作用，商由 \ref{prop:V.1.9} 为 $Y_1$. 设 $y_1$ 是 $Y_1$ 在 $y$ 之上的惟一点，由于 $k(y) = k(y_1)$，得出 $X$ 在 $y$ 处的纤维同构于 $X_1$ 在 $y$ 处的纤维，故 $X_1$ 在 $x$ 之上有惟一点 $x_1$. 此外由 \ref{prop:V.1.9} 有 $X_1/H = X'_1 = X' \times_Y Y_1$，设 $x'_1$ 是 $x_1$ 在 $X'_1$ 中的像，它在 $x'$ 之上，且容易验证 ($X'$ 在 $Y$ 上有限型) 同态 $\theta_{x'} \to \theta_{x'_1}$ 在完备化上诱导同构. 故归结到 $Y$ 是完备局部环 $B$ 的谱的情形，从而 $X$ 是 $B$ 上有限环 $A$ 的谱，它由对应于 $Y$ 上 $X$ 的诸点 $x_i$ 的有限个局部环 $A_x$ 组成. 若 $A_0$ 对应于 $x = x_0$，则 $A$ 等同于环 $\mathrm{Hom}_{\mathcal{G}_d}(G, A_0)$，即满足 $f(st) = sf(t)$ ($s \in G_d$) 的函数 $f \colon G \to A_0$ 的环，$G$ 在这些函数上的作用由 $(uf)(t) = f(tu)$ 定义. 于是看到，若 $H$ 是 $G$ 的任意子群，则 $A^{H}$ 是满足
\[
f(stu) = sf(t) \quad \text{对 } s \in G_d, u \in H
\]
的函数 $f \colon G \to A_0$ 的环，故它是半局部环，其局部分量对应于 $G$ 中的双陪集 $G_d a H$，对应于 $a \in G$ 的双陪集 (借助映射 $f \mapsto f(a)$) 的是 $A_0$ 的子环 $A_0^{H(a)}$，其中 $H(a) = G_d \cap aHa^{-1}$. 此外，$A^{H}$ 中对应于 $x$ 的像 $x'$ 的局部分量，也是对应于单位元素的双陪集 $G_d H$ 的那个分量，故其局部环为 $A_0^{G_d \cap H}$. 因此若 $G_d \subset H$，则得到 $A_0^{G_d} = A^{G} = B$，这证明 (i). 为证明 (ii)，可以对 $A$ 作适当的有限平坦扩张，用 \ref{prop:V.2.1}，归结到剩余扩张 $k(x)/k(y)$ 平凡的情形. 但此时 $G_i(x) = G_d(x)$，归结到前一情形.
\end{proof}

\begin{corollary}\label{cor:V.2.3}
在 \ref{prop:V.2.2} 的条件下，假设 $G_i(x) = (e)$，则 $X$ 在 $x$ 处平展于 $Y$ 上. 故若对每个 $x \in X$ 有 $G_i(x) = (e)$，则 $X \to Y$ 是平展态射.
\end{corollary}

有部分逆命题：

\begin{corollary}\label{cor:V.2.4}
假设 $X$ 连通且群 $G$ 在 $X$ 上忠实. 为使 $p \colon X \to Y = X/G$ 平展，必要且充分条件是 $X$ 的诸点的惯性群约化为单位元. 若如此，则 $G$ 等同于 $Y$-概形 $X$ 的所有 $Y$-自同构的群.
\end{corollary}

\begin{proof}
考虑到 \ref{cor:V.2.3}，可以假设 $X$ 在 $Y$ 上平展. 但若 $s \in G$ 属于某个 $G_i(x)$，则由 \ref{cor:I.5.4} 得出 $s$ 在 $G$ 上平凡作用，故是单位元，因为 $G$ 忠实，这证明第一个断言. 设 $u$ 是 $X$ 的 $Y$-自同构，设 $x \in X$. 由命题 \ref{prop:V.1.3}，存在 $s \in G$，使得 $s(x) = u(x)$，且诱导与 $u$ 相同的剩余同态 $k(x)/k(x')$. 由上引文献 (loc. cit.) 有 $s = u$，证毕.
\end{proof}

\begin{remark}\label{rem:V.2.5}
在推论 \ref{cor:V.2.4} 中``$G$ 忠实''的假设当然不是多余的. ``$X$ 连通''的假设也不是，例如取 $X = Y \times E$，$E$ 是有限集，$G$ 是 $E$ 的置换群：$G$ 的作用带强惯性，然而 $(Y \times E)/G = Y \times (E/G) = Y$，且 $X$ 在 $Y$ 上平展. 取 $G$ 为严格小于 $E$ 的对称群、但传递作用在 $E$ 上的群，看到还有不来自 $G$ 的 $X$ 的 $Y$-自同构.
\end{remark}

无惯性作用的典型例子是 $Y \times G$，让 $G$ 借助它在因子 $G$ 上的右平移作用：右算子群为 $G$ 的 $Y$-预概形 $X$ 称为平凡的，若它同构于 $Y \times G$.

为建立带有限算子群的预概形与范畴中主丛 (principal bundle) 概念之间的联系 (这一联系对本讨论班其余部分虽不需要，但在其他语境中很重要)，下面的考虑是有用的. 固定基预概形 $Y$，处于 $Y$-预概形范畴中. 若 $G$ 是有限群，为简短令 $G_Y = Y \times G$，它是 $Y$ 上的有限群概形 (参见第 1 节)，若 $X$ 是 $Y$-预概形，有
\[
X \times_Y G_Y = X \times G
\]
(同一参考). $Y$-态射 $X \times_Y G_Y \to X$ 的给出因此等价于 $Y$-态射 $X \times G \to X$ 的给出，i.e. 对每个 $g \in G$ 给出 $Y$-态射 $T_g \colon X \to X$. 立即看到，为使诸 $T_g$ 在 $X$ 上定义右算子群 $G$ 的预概形结构 (i.e. $T_{gg'} = T_{g'}T_g$，$T_e = \mathrm{id}_x$)，必要且充分条件是相应的 $Y$-态射 $X \times_Y G_Y \to X$ 在 $X$ 上定义 $Y$-群概形 $Y$-算子群的 $Y$-预概形结构 (在范畴 $\mathcal{C}$ 中 $\mathcal{C}$-算子群对象的普遍意义下). 假设如此. 回忆若典范态射
\[
X \times_Y G_Y \longrightarrow X \times_Y X
\]
(其分量分别为 $\mathrm{pr}_1$ 和乘法态射 $\pi \colon X \times_Y G_Y \to X$) 是同构，则称 $X$ 在 $G_Y$ 下形式主齐性\footnote{现在宁愿说：$X$ 是 $G_Y$ 下的伪挠子 (pseudo-torseur sous $G_Y$).}. 在目前情形，把第一个成员等同于 $X \times G$，所考虑的态射是给每个 $g \in G$ 对应态射
\[
(\operatorname{id}_X, T_g) = (\operatorname{id}_X \times_Y T_g)\Delta_{X/Y} \colon X \longrightarrow X \times_Y X
\]
的那个，从而说 $X$ 在 $G_Y$ 下形式主齐性，也意味着 $X \times_Y X$ 同构于对角线在 $G \times G$ 的元素 $(e, g)$ 下的变换的直和 ($G \times G$ 以显然的方式作用在 $X \times_Y X$ 上)，其中 $e$ 是 $G$ 的单位元. 若不想区分左右、给出对多于两个等同于 $X$ 的因子乘积仍适用的公式，可以把条件表述为：典范态射
\[
X \times_G (G \times G) \longrightarrow X \times_Y X
\]
(它由给配对 $(g, g')$ 对应态射
\[
(T_g, T_{g'}) = (T_g \times_Y T_{g'})\Delta_{X/Y} \colon X \longrightarrow X \times_Y X
\]
并使 $G$ 通过对角同态左作用在 $G \times G$ 上：
\[
s(g, g') = (sg, sg')
\]
得到) 是同构.

主齐性空间 (principal homogeneous space) 的概念由形式主齐性空间的概念添加一个公理得出，该公理保证 $X$ 除以 $G_Y$ 的``商''存在，且正是范畴的右幺元对象，这里即 $Y$. 这个公理依语境可以变化，在``下降''的框架中通常最方便地明确为：要求算子对象经一个合适类型的合适基变换后成为``平凡''的，i.e. 同构于乘积 $X \times_Y G_Y$ (此处即 $X \times G$)(如此，实际上允许下降技术；参见 Grothendieck，Technique de descente et théorèmes d'existence en Géométrie Algébrique，Sém. Bourbaki N$^{\circ}$ 190，第 26 至 28 页)\footnote{关于平坦下降理论参见讲 VIII.}. 在这种思路下，这里指出 $G$ 群主齐性丛 (在上引文献 (loc. cit.) 的意义下) 的刻画：

\begin{proposition}\label{prop:V.2.6}
设 $Y$ 是局部诺特预概形，$X$ 是带右作用的有限算子群 $G$ 的 $Y$-预概形. 下列条件等价：
\begin{enumerate}[(i)]
\item $X$ 在 $Y$ 上有限，$Y = X/G$，$X$ 的诸点的惯性群约化为单位元；
\item 存在忠实平坦拟紧的基变换 $Y_1 \to Y$，使得 $X_1 = X \times_Y Y_1$ 是平凡的 $Y_1$-算子预概形，i.e. 同构于 $Y_1 \times G$；
\item[(ii bis)] 同 (ii)，但 $Y_1 \to Y$ 有限、平展、满射；
\item $X$ 在 $G_Y$ 下形式主齐性，且在 $Y$ 上忠实平坦且拟紧.
\end{enumerate}
\end{proposition}

\paragraph*{证明.} (i)$\Rightarrow$(ii bis). 取 $Y_1 = X$，注意 $X \to Y$ 确为有限、平展 (由 \ref{cor:V.2.3}) 且满射. 证明 $X_1$ 于是在 $Y_1$ 上平凡，这由下述推论得出：

\begin{corollary}\label{cor:V.2.7}
若 (i) 成立且 $X$ 在 $Y$ 上有截面，则 $X$ 是平凡算子空间.
\end{corollary}

\begin{proof}
事实上，这个截面可以定义 $G$-态射 $X \times G \to X$，由于 $G$ 在 $X$ 的纤维上传递它满射，由于 $G$ 无惯性作用它单射；最后，由 \ref{cor:I.5.3}，因为 $X$ 在 $Y$ 上平展，它是局部同构. 故它是同构.
\end{proof}

(ii bis) 平凡蕴含 (ii)，(ii) 蕴含 (i)，因为 (i) 的成分对基的忠实平坦拟紧扩张``不变''(关于``有限''参见上面引用的 Bourbaki 讨论班；对惯性群应用 \ref{prop:V.2.1}；对 $Y = X/G$，用 \ref{prop:V.1.9} 在忠实平坦基变换情形的逆命题，我们忘了把它明确化).

我们通过证明 (i)$\Rightarrow$(ii bis) 证明了 (i)$\Rightarrow$(iii). 最后 (iii)$\Rightarrow$(ii)，因为 (iii) 中的第一个假设恰意味着 $X$ 在作基变换 $Y_1 = X$ 后平凡；由此 (ii)，因为 $X$ 在 $Y$ 上忠实平坦且拟紧.

\begin{definition}\label{def:V.2.8}
满足 \ref{prop:V.2.6} 等价条件的带右算子群 $G$ 的 $Y$-预概形 $X$，称为 $Y$ 的以 $G$ 为 伽罗瓦群的主覆盖 (principal covering).
\end{definition}

\section{平展覆盖的自同构与态射}

\begin{proposition}\label{prop:V.3.1}
设 $X$ 在局部诺特 $Y$ 上平展、分离、有限型，$G$ 是通过 $Y$-自同构作用在 $X$ 上的有限群. 则 $G$ 的作用可允许，且商预概形 $X/G$ 在 $Y$ 上平展.
\end{proposition}

\begin{proof}
不假设 $X$ 在 $Y$ 上有限，但 $X$ 在 $Y$ 上拟射影，从而由 \ref{cor:V.1.8} 存在 $X/G$. 先证明：
\end{proof}

\begin{corollary}\label{cor:V.3.2}
态射 $X \to X/G$ 平展.
\end{corollary}

\begin{proof}
显然可以先假设 $G$ 在 $X$ 的连通分支集合上传递，然后考虑一个连通分支的稳定化子，甚至可以假设 $X$ 连通. 最后可以假设 $G$ 忠实作用. 但此时像 \ref{cor:V.2.4} 中那样看到 $G$ 无惯性作用，故由 \ref{cor:V.2.3} 得出 $X \to X/G$ 平展. 借助下述引理得出结论：
\end{proof}

 \begin{lemma}[对讲 I 的更正]\label{lem:V.3.3}
设 $X \to X' \to Y$ 是有限型态射，$x$ 是 $X$ 的点，$x'$ 和 $y$ 是它的像. 假设 $Y$ 局部诺特. 若所考虑的三个态射中有两个在标记点处平展，则第三个亦然.
\end{lemma}

\begin{proof}
只剩下 $X \to X'$ 和 $X \to Y$ 在 $x$ 处平展、要证明 $X' \to Y$ 在 $x'$ 处平展的情形 (这正是 \ref{prop:V.3.1} 所需). 作 $Y$ 的适当平坦扩张，归结到剩余扩张 $k(x)/k(y)$ 平凡的情形. 考虑同态 $\mathcal{O}_y \to \mathcal{O}_{x'} \to \mathcal{O}_x$ 及由完备化得到的同态，假设意味着 $\widehat{\theta_y} \to \widehat{\theta_x}$ 和 $\widehat{\mathcal{O}_{x'}} \to \widehat{\mathcal{O}_x}$ 是同构，由此立即得出 $\widehat{\mathcal{O}_y} \to \widehat{\mathcal{O}_{x'}}$ 是同构，这证明引理.
\end{proof}

\begin{corollary}\label{cor:V.3.4}
若 $X$ 在 $Y$ 上有限且平展，则 $X/G$ 在 $Y$ 上有限且平展.
\end{corollary}

\begin{proposition}\label{prop:V.3.5}
设 $X$，$X'$ 是 $Y$ 的两个平展覆盖. 则每个 $Y$-态射 $f \colon X \to X'$ 分解为平展满射态射 $X \to X''$ 与 $X'$ 的一个既开又闭子集 $X''$ 的典范浸入 $X'' \to X'$ 的乘积.
\end{proposition}

\begin{proof}
知道 (\ref{cor:I.4.8}) $f$ 平展，故开；另一方面 $X$ 在 $Y$ 上有限，故 $f$ 闭，因此 $f(X) = X''$ 是 $X'$ 的既开又闭子集. 证明完毕 (N.B. 只需 $X'$ 在 $Y$ 上非分歧，而不必是平展覆盖).
\end{proof}

\begin{corollary}\label{cor:V.3.6}
用前面的记号，$X \to X''$ 是预概形范畴中的严格满态射 (strict epimorphism)，$X'' \to X'$ 是预概形范畴中的单态射 (monomorphism) (甚至是严格单态射).
\end{corollary}

\begin{proof}
第一个断言按定义意味着态射序列
\[
X \times_{X''} X \xrightarrow[\mathrm{pr}_2]{\mathrm{pr}_1} X \longrightarrow X''
\]
正合，这由 $X \to X''$ 有限且忠实平坦的事实得出 (容易看出，参见 Grothendieck，loc. cit.). 对 $X'' \to X'$ 的对偶断言更平凡.
\end{proof}

推论 \ref{cor:V.3.6} 在下一节的基本群理论中对我们有用；(对那些不喜欢严格满态射概念的人) 可以把推论 \ref{cor:V.3.6} 换成读者按个人口味安排的变体. 只借此机会指出：分解 $f = f'f''$ (其中 $f''$ 是严格满态射，$f'$ 是单态射) 在 (任何范畴中) 唯一同构意义下必然唯一；然而，可以同时存在具有对偶性质的分解 $f = f_1 f_2$ ($f_2$ 是满态射，$f_1$ 是严格单态射，也在唯一同构意义下唯一)，它不同构于前一个：只需取例如拓扑向量空间 (若愿意，分离的) 的范畴，及使 $u(X)$ 不闭的 $u \colon X \to X'$.

\begin{proposition}\label{prop:V.3.7}
设 $Y$ 是连通局部诺特预概形，$y$ 是 $Y$ 的点，$\Omega$ 是 $k(y)$ 的代数闭扩张. 对 $Y$ 上的每个 $X$，用 $X(\Omega)$ 表示 $X$ 在 $\Omega$ 中取值的点的集合. 设 $X$，$X'$ 是 $Y$ 的平展覆盖，$u \colon X \to X'$ 是 $Y$-态射，使得相应的映射 $X(\Omega) \to X'(\Omega)$ 是同构. 则 $u$ 是同构.
\end{proposition}

\begin{proof}
立即归结到 $X'$ 连通的情形. 由于 $X \to X'$ 有限且平展，知道 $X \to X'$ 的纤维中点的几何个数恒定，且等于 1 当且仅当所考虑的态射是同构. 而这个数也是 $X(\Omega) \to X'(\Omega)$ 的一个纤维中的元素个数，由此得结论.
\end{proof}

\section{Galois 理论的公理条件}

设 $\mathcal{C}$ 是范畴，$F$ 是 $\mathcal{C}$ 到有限集范畴的共变函子. 假设下列条件满足：
\begin{itemize}\tightlist
\item[(G 1)] $\mathcal{C}$ 有终对象 (final object)\footnote{回忆 $\mathcal{C}$ 的对象 $e$ 称为终对象，若对 $\mathcal{C}$ 中的每个 $X$，$\mathrm{Hom}(X, e)$ 恰好有一个元素. 对偶地定义 $\mathcal{C}$ 的始对象 (initial object).}，且 $\mathcal{C}$ 中两个对象在第三个之上的纤维积存在 (这个公理也可以说成：$\mathcal{C}$ 中存在有限投射极限).
\item[(G 2)] $\mathcal{C}$ 中存在有限直和 (故也存在扮演空集角色的始对象 $\varnothing_{\mathcal{C}}$)，以及对象除以有限自同构群的商.
\item[(G 3)] 设 $u \colon X \to Y$ 是 $\mathcal{C}$ 中的态射，则 $u$ 分解为乘积 $X \xrightarrow{u'} Y' \xrightarrow{u''} Y$，其中 $u'$ 是严格满态射，$u''$ 是单态射，且是 $Y$ 的一个直因子上的同构.
\item[(G 4)] 函子 $F$ 左正合 (i.e. 把右幺元变为右幺元，且与纤维积可交换).
\item[(G 5)] $F$ 与有限直和可交换，把严格满态射变为满态射，且与除以有限自同构群的商可交换.
\item[(G 6)] 设 $u \colon X \to Y$ 是 $\mathcal{C}$ 中的态射，使得 $F(u)$ 是同构，则 $u$ 是同构.
\end{itemize}

我们的目标是从条件 (G 1) 至 (G 6) 构造一个投射有限 (projective limit of finite groups) 的拓扑群 $\pi$，以及 $\mathcal{C}$ 与有限集范畴 $\mathcal{C}(\pi)$\index[notations]{$\mathcal{C}(\pi)$} (其中 $\pi$ 连续作用，i.e. 作用的稳定化子是开子群，或等价地，已作用在所考虑集合上的一个离散商群存在) 的等价，所构造的等价把给定函子 $F$ 变为 $\mathcal{C}(\pi)$ 到有限集范畴的显然包含函子. 立即注意，用拓扑群 $\pi$ 构造的范畴 $\mathcal{C}(\pi)$ 和前面的包含函子确实满足条件 (G 1) 至 (G 6).

分几步进行.

\begin{enumerate}[a)]
\item 设 $u \colon X \to Y$ 在 $\mathcal{C}$ 中. 为使 $u$ 是单态射，必要且充分条件是 $F(u)$ 单射. (用 (G 1)，(G 4)，(G 6).)
\end{enumerate}

事实上，说 $u$ 是单态射意味着投影 $\mathrm{pr}_1 \colon X \times_Y X \to X$ 是同构.

\begin{enumerate}[a)]
\setcounter{enumi}{1}
\item 对象 $X$ 都是 $\mathcal{C}$ 的 Artin 对象.
\end{enumerate}

事实上，若 $X' \to X'' \to X$ 是单态射，使得 $F(X')$ 和 $F(X'')$ 在 $F(X)$ 中有相同的像，则由 a) $F(X') \to F(X'')$ 是同构，故由 (G 6)，$X' \to X''$ 是同构.

\begin{enumerate}[a)]
\setcounter{enumi}{2}
\item 函子 $F$ 严格 pro-可表示 (strictly pro-representable). (参见 Grothendieck，Technique de descente et théorèmes d'existence en Géométrie Algébrique，II，Séminaire Bourbaki 195，février 1960.)
\end{enumerate}

事实上，由上引文献 (loc. cit.) 命题 \ref{prop:V.3.1}，这由 b) 和 (G 4) 得出. 故可以找到 $I$ 上的有序滤过投射系：
\[
P = (P_i)_{i \in I}
\]
(看作 $\mathcal{C}$ 的 pro-对象)，以及函子同构
\begin{equation}
F(X) = \operatorname{Hom}_{\operatorname{Pro}(\mathcal{C})}(P, X) = \varinjlim_{i} \operatorname{Hom}_{\mathcal{C}}(P_i, X). \tag{($*$)}
\end{equation}
该同构由元素
\[
\varphi \in \varprojlim_{i} F(P_i) = F(P)
\]
实现. 还可以假设转移同态 $\varphi_{ji} \colon P_i \to P_j$ ($i \geqslant j$) 是满态射，且每个满态射 $P_i \to P'$ 等价于对适当的 $j \leqslant i$ 的满态射 $P_i \to P_j$ (这基本上唯一地确定投射系 $P$).

称 $X \in \mathcal{C}$ 连通 (connected)，若它不同构于两个不都同构于始对象 $\varnothing_{\mathcal{C}}$ 的对象的和.

\begin{enumerate}[a)]
\setcounter{enumi}{3}
\item 诸 $P_i$ 连通且不同构于 $\varnothing_{\mathcal{C}}$.
\end{enumerate}

若 $X$ 是左幺元，由应用于空族的直和的 (G 5)，$F(X) = \varnothing$，反之亦然 (由 (G 6)). 故若 $X'$ 是 $\mathcal{C}$ 中非左幺元的对象，i.e. $F(X') \neq \varnothing$，则不存在从 $X'$ 到 $X$ 的态射. 故若某个 $P_i$ 是左幺元，则 $i$ 是有序滤过指标集 $I$ 的最大元素，且公式 (∗) 将意味着 $F(X) = \mathrm{Hom}(P_i, X) =$ 对每个 $X$ 约化为一个元素的集合，这荒谬，因为 $F(\varnothing_{\mathcal{C}}) = \varnothing$. 故诸 $P_i$ 不同构于 $\varnothing_{\mathcal{C}}$.

假设 $P_i = A \amalg B$，由此由 (G 5) $F(P_i) = F(A) \amalg F(B)$，特别是 $F(P_i)$ 中对应于 (由 (∗)) 恒等同态 $P_i \to P_i$ 的元素 $a_i$ 在 $F(A) \amalg F(B)$ 中，例如在 $F(A)$ 中. 这意味着存在 $j \geqslant i$，使得 $\varphi_{ij} \colon P_j \to P_i$ 分解为 $P_j \to A \to P_i = A \amalg B$，其中第二个箭头是典范态射. 故 $F(P_j) \to F(P_i)$ 分解为 $F(P_j) \to F(A) \to F(P_i) = F(A) \amalg F(B)$，且由于由 (G 5)，$F(P_j) \longrightarrow F(P_i)$ 满射，得出 $F(B) = \varnothing$，故 $B$ 同构于 $\varnothing_{\mathcal{C}}$.

\begin{enumerate}[a)]
\setcounter{enumi}{4}
\item 每个态射 $u \colon X \to Y$ (其中 $X$ 不同构于 $\varnothing_{\mathcal{C}}$，$Y$ 连通) 都是严格满态射. 连通对象的每个自同态都是自同构.
\end{enumerate}

考虑 $u$ 的 (G 3) 分解，由于 $X \neq \varnothing_{\mathcal{C}}$，由 (G 6) 得出 $F(X) \neq \varnothing$，故 $F(Y') \neq \varnothing$，从而 $Y' \neq \varnothing_{\mathcal{C}}$，故 $Y$ 连通时 $Y'$ 等同于 $Y$，因此 $u$ 是严格满态射. 假设 $u$ 是连通对象 $X$ 的自同态，证明它是自同构. 事实上，可以假设 $X$ 不同构于 $\varnothing_{\mathcal{C}}$，故由前面所述 $u$ 是严格满态射，因此由 (G 5) $F(u)$ 是满态射，又因 $F(X)$ 是有限集，$F(u)$ 双射. 故由 (G 6)，$u$ 是自同构.

特别地，每个 $P_i$ 的自同态都是自同构.

\begin{enumerate}[a)]
\setcounter{enumi}{5}
\item 对 $P_i$ 的下列条件等价：(i) 自然单射 $\mathrm{Hom}(P_i, P_i) \to \mathrm{Hom}(P, P_i) \simeq F(P_i)$ 也满射，i.e. 对每个 $u \colon P \to P_i$，存在 $v \colon P_i \to P_i$ 使得 $u = v\varphi_i$ (其中 $\varphi_i$ 是典范同态 $P \to P_i$). (ii) 群 $\operatorname{Aut}(P_i)$ 传递作用在 $F(P_i)$ 上. (iii) 群 $\operatorname{Aut}(P_i)$ 简单传递作用在 $F(P_i)$ 上.
\end{enumerate}

事实上，把 $\mathrm{Hom}(P, P_i)$ 等同于 $F(P_i)$，(i) 中所考虑的映射正是 $v \mapsto F(v)(\varphi_i)$. 三个条件的等价性来自 $\mathrm{Hom}(P_i, P_i) = \operatorname{Aut}(P_i)$ 且前面的映射已经单射这一事实.

满足 f) 的等价条件 (i)，(ii)，(iii) 的 $P_i$ 称为 Galois 的 (Galois).

\begin{enumerate}[a)]
\setcounter{enumi}{6}
\item 对 $\mathcal{C}$ 中的每个 $X$，存在 Galois 的 $P_i$，使得每个 $u \in \mathrm{Hom}(P, X)$ 分解为 $P \xrightarrow{\varphi_i} P_i \longrightarrow X$.
\end{enumerate}

设 $J = \operatorname{Hom}(P, X) = F(X)$，它是有限集，故存在 $P_j$，使得每个 $u \colon P \to X$ 分解为 $P \xrightarrow{\varphi_j} P_j \longrightarrow X$，或等价地，使得自然态射
\[
P \longrightarrow X^{J} \quad (J = \operatorname{Hom}(P, X))
\]
分解为
\[
P \xrightarrow{\varphi_j} P_j \longrightarrow X^{J}.
\]
由 (G 3)，态射 $P_j \to X^{J}$ 分解为单态射与严格满态射的乘积，可以取为形如 $\varphi_{ij} \colon P_j \to P_i$. 故归结为证明 $P_i$ Galois. 设 $k \geqslant j$ 是使得每个态射 $P \to P_i$ 都经 $P \stackrel{\varphi_k}{\longrightarrow} P_k \longrightarrow P_i$ 分解的指标. 注意自然态射 $P_k \to X^{J}$ 仍分解为复合
\[
P_k \xrightarrow{\varphi_{ik}} P_i \xrightarrow{U} X^{J},
\]
其中第一个箭头由 e) 是严格满态射，第二个是单态射. 要证明：对给定的态射 $\psi \colon P_k \to P_i$，存在 $P_i$ 的自同态 $v$ 使得 $\psi = v\varphi_{ik}$. 但对每个 $u \in \mathrm{Hom}(P_i, X)$，考虑 $u\psi \in \mathrm{Hom}(P_k, X)$，它形如 $u'\varphi_{ik}$，其中 $u' \in \mathrm{Hom}(P_i, X)$ 完全确定. 由 $\psi$ 如此定义的 $J$ 到 $J$ 的映射 $u \mapsto u'$ 还单射，因为由 e) $\psi$ 是满态射；由于 $J$ 有限它双射. 故 $J$ 到 $J$ 的双射映射 $u \mapsto u'$ 定义同构 $\alpha \colon X^{J} \xrightarrow{\sim} X^{J}$，使下面的图交换. 由态射分解为单态射与严格满态射的乘积的唯一性，得出 (由于 $\psi$ 由 e) 也是严格满态射) 可以找到使图交换的态射 $v \colon P_i \to P_i$，证毕.

特别得出，Galois 的 $P_i$ 在诸 $(P_j)$ 中构成共尾系. 由于对 Galois 对象 $P_i$ 有
\[
\operatorname{Hom}(P, P_i) = \operatorname{Hom}(P_i, P_i) = \operatorname{Aut}(P_i),
\]
取极限得：
\[
\operatorname{Hom}(P, P) = \varprojlim_{i} \operatorname{Hom}(P, P_i) = \varprojlim_{i} \operatorname{Hom}(P_i, P_i) = \varprojlim_{i} \operatorname{Aut}(P_i),
\]
其中投射极限在 Galois 的 $P_i$ 上取. 此外，借助等同 $\mathrm{Hom}(P, P_i) = F(P_i)$，并考虑到 $F$ 把满态射变为满态射，看到前面投射系的转移同态满射. 由这一切得出结论：

\begin{enumerate}[a)]
\setcounter{enumi}{7}
\item 有
\[
\operatorname{Hom}(P, P) = \operatorname{Aut}(P) = \varprojlim_{i} F(P_i) = \varprojlim_{i} \operatorname{Aut}(P_i),
\]
其中投射极限在 Galois 的 $P_i$ 上取.
\end{enumerate}

特别地，$\operatorname{Aut}(P)$ 表现为有限群投射系 (转移同态满射) 的投射极限，给它配备离散拓扑的投射极限拓扑. 记 $\pi$ 并称为\index[notations]{$\pi$} ($\mathcal{C}$ 配备 $F$ 的) \emph{基本群} (fundamental group)\index[term]{jibenqun@基本群} 的是 $\operatorname{Aut}(P)$ 的对偶群 (opposite group). 该群因此右作用在 $P$ 上，它是右作用在 Galois 的 $P_i$ 上的有限群 $\pi_i$ 的投射极限，$\pi_i$ 是 $\operatorname{Aut}(P_i)$ 的对偶群.

考虑到函子同构
\[
F(X) = \operatorname{Hom}(P, X)
\]
和 $\pi$ 的定义，看到 $\pi$ 左作用在 $F(X)$ 上，且由 g) 连续作用 (因为用 g) 的记号，实际上是 $\pi_i$ 作用在 $F(X)$ 上). 对 $\mathcal{C}$ 中的每个态射 $u \colon X \to Y$，态射 $F(u) \colon F(X) \to F(Y)$ 与 $\pi$ 的作用相容，这是平凡的. 因此可以把 $F$ 看作共变函子
\[
F \colon \mathcal{C} \longrightarrow \mathcal{C}(\pi),
\]
其中 $\mathcal{C}(\pi)$ 是 $\pi$ 在其上连续左作用的有限集的范畴.

现在定义反方向的函子：
\[
G \colon \mathcal{C} \longleftarrow \mathcal{C}(\pi)
\]
由公式
\[
G(E) = P \times_\pi E,
\]
其中 $P \times_\pi E$ 定义为由下式概括的万有问题的解
\[
\operatorname{Hom}_{\mathcal{C}}(P \times_\pi E, X) \xrightarrow{\sim} \operatorname{Hom}_\pi(E, \operatorname{Hom}(P, X))
\]
(其中第二个成员中的 $\mathrm{Hom}(P, X) = F(X)$ 看作 $\pi$ 在其上左作用的集合). 必须证明对象 $P \times_\pi E$ 的存在性.

\begin{enumerate}[a)]
\setcounter{enumi}{8}
\item 设 $Q$ 是有限群 $G$ 在其上右作用的 $\mathcal{C}$ 的对象，$E$ 是 $G$ 在其上左作用的有限集. 则 $Q \times_G E$ 存在，且典范映射
\[
F(Q) \times_G E \longrightarrow F(Q \times_G E)
\]
是同构.
\end{enumerate}

由于由 (G 2) $\mathcal{C}$ 中存在有限直和，且 $F$ 由 (G 5) 与其可交换，立即归结到 $G$ 传递作用在 $E$ 上的情形，因为若 $E_j$ 是 $G$ 在 $E$ 中的轨道，则有
\[
Q \times_G E = \coprod_{j} Q \times_G E_j.
\]
设 $a \in E$，$H$ 是其稳定化子，由定义立即看到 $Q \times_G E$ 等同于 $Q/H$. 由此得到存在性 (借助 (G 2))，以及对 $F$ 的交换性 (借助 (G 5)).

\begin{enumerate}[a)]
\setcounter{enumi}{9}
\item 设 $E$ 是 $\mathcal{C}(\pi)$ 的对象，$P_i$ 是 $\pi_i$ 已作用在 $E$ 上的 Galois 对象. 则 $P_i \times_{\pi_i} E$ 存在，且有典范同构
\[
E \stackrel{\sim}{\longrightarrow} F(P_i \times_{\pi_i} E).
\]
若 $j \geqslant i$ 使得 $P_j$ Galois，则典范同态 $P_j \times_{\pi_j} E \to P_i \times_{\pi_i} E$ 是同构.
\end{enumerate}

第一个断言是 i) 的特殊情形，考虑到 $\pi_i$ 简单传递作用在配备标记点 $\varphi_i$ 的 $F(P_i)$ 上，由此得同构 $F(P_i) \times_{\pi_i} E \simeq E$. 对第二个断言例如用 (G 6).

对每个 $j$，设 $\mathcal{C}_j$ 是 $\mathcal{C}$ 的由满足``$\mathrm{Hom}(P_j, X) \to \mathrm{Hom}(P, X) \simeq F(X)$ 双射''的 $X$ 组成的全子范畴. 由 g) 知道 $\mathcal{C}$ 是诸 $\mathcal{C}_j$ 的滤过并. 故对 $X \in \mathcal{C}_j$ 有
\[
\begin{array}{r}
\mathrm{Hom}_\pi(E, \mathrm{Hom}(P, X)) \simeq \mathrm{Hom}_\pi(E, \mathrm{Hom}(P_j, X)) \simeq \mathrm{Hom}_{\pi_j}(E, \mathrm{Hom}(P_j, X)) \\
\simeq \mathrm{Hom}(P_j \times_{\pi_j} E, X),
\end{array}
\]
考虑到 j) 中最后一个断言，得到对 $\mathcal{C}_j$ 的对象 $X$ 函子性的同构：
\[
\mathrm{Hom}_\pi(E, \mathrm{Hom}(P, X)) \simeq \mathrm{Hom}(P_i \times_{\pi_i} E, X).
\]
由于这对每个 $j$ 成立，且这些函子同构随 $j$ 变动而互相诱导，得出结论：

\begin{enumerate}[a)]
\setcounter{enumi}{10}
\item 在 j) 的条件下，由典范态射复合而成的态射
\[
E \longrightarrow \operatorname{Hom}(P_i, P_i \times_{\pi_i} E) \longrightarrow \operatorname{Hom}(P, P_i \times_{\pi_i} E)
\]
使 $P_i \times_{\pi_i} E$ 成为定义 $P \times_\pi E$ 的万有问题的解，i.e. 后者存在且有同构
\[
P \times_\pi E \stackrel{\sim}{\longrightarrow} P_i \times_{\pi_i} E.
\]
\end{enumerate}

这完成函子 $G(E)$ 的构造. 另一方面有函子同态
\[
\alpha \colon \mathrm{id}_{\mathcal{C}(\pi)} \longrightarrow FG,
\]
i.e. 对 $\mathcal{C}(\pi)$ 的对象 $E$ 函子性的同态：
\[
\alpha(E) \colon E \longrightarrow FG(E) = F(P \times_\pi E),
\]
即典范态射的复合
\[
E \longrightarrow F(P) \times_\pi E \longrightarrow F(P \times_\pi E)
\]
(其中第一个来自标记点 $\varphi \in F(P)$). 结合 j) 和 k)，得到：

\begin{enumerate}[a)]
\setcounter{enumi}{11}
\item 同态 $\alpha$ 是同构.
\end{enumerate}

类似地定义函子同态
\[
\beta \colon GF \longrightarrow \mathrm{id}_{\mathcal{C}},
\]
i.e. 对 $\mathcal{C}$ 的对象 $X$ 函子性的同态：
\[
\beta(X) \colon P \times_\pi F(X) \longrightarrow X,
\]
作为与 $\pi$-同态
\[
F(X) \longrightarrow \operatorname{Hom}(P, X)
\]
(它是典范同构 $\mathrm{Hom}(P, X) \stackrel{\sim}{\longrightarrow} F(X)$ 的逆) 相伴的同态.

\begin{enumerate}[a)]
\setcounter{enumi}{12}
\item 复合
\[
\begin{array}{c}
F(X) \xrightarrow{\alpha(F(X))} FGF(X) \xrightarrow{F(\beta(X))} F(X), \\
G(E) \xrightarrow{G(\alpha(E))} GFG(E) \xrightarrow{\beta(G(E))} G(E)
\end{array}
\]
是恒等同构.
\end{enumerate}

(此处原书有一句无法辨认的旁注.)

考虑到 l) 得出：

\begin{enumerate}[a)]
\setcounter{enumi}{13}
\item 同态 $\beta$ 是同构.
\end{enumerate}

这样得到所承诺的结果：

\begin{theorem}\label{thm:V.4.1}
设 $\mathcal{C}$ 是满足本节开头条件 (G 1)，(G 2)，(G 3) 的范畴，$F$ 是 $\mathcal{C}$ 到有限集范畴的共变函子，满足条件 (G 4)，(G 5)，(G 6). 则前面的典范构造定义互为拟逆的\emph{范畴等价} $F \colon \mathcal{C} \to \mathcal{C}(\pi)$ 和 $G \colon \mathcal{C}(\pi) \to \mathcal{C}$. 精确地说，存在 $\mathcal{C}$ 的 pro-对象 $P$ 和函子同构 $F(X) \longleftarrow \mathrm{Hom}(P, X)$，$\pi$ 是 $P$ 的自同构群的对偶群，以合适的方式拓扑化，使得 $\pi$ 连续作用在集合 $\mathrm{Hom}(P, X) \simeq F(X)$ 上. 最后，$G$ 由 $G(E) \simeq P \times_\pi E$ 给出.
\end{theorem}

\begin{remark}\label{rem:V.4.2}
若把 (G 2) 和 (G 5) 分别换成：
\begin{itemize}\tightlist
\item[(G' 2)] $\mathcal{C}$ 中存在有限归纳极限；
\item[(G' 5)] 函子 $F$ 右正合 (i.e. 与有限归纳极限可交换)，
\end{itemize}
则条件 (G 1) 至 (G 6) 的陈述变得更简单更友好. 这些条件表面上强于 (G 2) 和 (G 5)，但由结构定理 \ref{thm:V.4.1} 立即得出它们被 (G 1) 至 (G 6) 蕴含. 然而注意，在我们关心的情形，(G 2) 和 (G 5) 的验证似乎确实比 (G' 2) 和 (G' 5) 简单. 我不知道在条件 (G 3) 中，``$u''$ 是 $Y$ 的一个直因子上的同构''这一事实是否可以被省略.
\end{remark}

\section{Galois 范畴}

\begin{definition}\label{def:V.5.1}
称范畴 $\mathcal{C}$ 为 \emph{Galois 范畴} (Galois category)\index[term]{Galois fanchou@Galois 范畴}，若它等价于某个范畴 $\mathcal{C}(\pi)$，其中 $\pi$ 是紧致群，有限群的投射极限 (i.e. 完全不连通).
\end{definition}

$\mathcal{C}(\pi)$ 的定义参见第 4 节开头. 由定理 \ref{thm:V.4.1}，$\mathcal{C}$ 是 Galois 范畴当且仅当它满足条件 (G 1) 至 (G 3)，且存在 $\mathcal{C}$ 到有限集范畴的、满足条件 (G 4) 至 (G 6) 的函子 $F$ (i.e. 在一般术语中，正合且保守的函子). 这样的函子称为 Galois 范畴 $\mathcal{C}$ 的\emph{基本函子} (fundamental functor)\index[term]{jibenhanzi@基本函子}\footnote{看来最好采用更形象的术语``\emph{纤维函子}'' (fiber functor).}；它被我们记为 $P_F$ 的 pro-对象 pro-表示；使得相伴函子 $F$ 基本 (fundamental) 的 pro-对象 $P$ 称为\emph{基本 pro-对象} (fundamental pro-object)\index[term]{jiben pro-duixiang@基本 pro-对象}. 这样，$\mathcal{C}$ 上的基本函子范畴与基本 pro-对象范畴反等价；若 $F$ 和 $P$ 对应，则群 $\mathrm{Aut}\, F$ 同构于群 $\mathrm{Aut}\, P$ 的对偶，故前一节记为 $\pi$ 的群正是 $\mathrm{Aut}\, P$. 回忆在前一节中我们从给定的基本函子 $F$ 构造了 $\mathcal{C}$ 与 $\mathcal{C}(\pi)$ (其中 $\pi = \operatorname{Aut}(F)$) 的等价，它把函子 $F$ 变为 $\mathcal{C}(\pi)$ 到有限集范畴的典范函子. 在这个典型情形 $\mathcal{C} = \mathcal{C}(\pi)$，$F =$ 典范函子中，与 $F$ 相伴的基本 pro-对象正是 $\pi$ 的离散商 $\pi_i$ 的投射系.

明确化 $\mathcal{C}(\pi)$ 的 pro-对象范畴可能是有用的. 得到：

\begin{proposition}\label{prop:V.5.2}
范畴 Pro-$\mathcal{C}(\pi)$ 典范等价于拓扑群 $\pi$ 的紧致完全不连通算子空间范畴 $\mathcal{C}'(\pi)$.
\end{proposition}

\begin{proof}
由于后者把 $\mathcal{C}(\pi)$ 作为全子范畴 (对应于紧致离散算子空间) 包含，且其中存在投射极限，无论如何有典范函子 $g \colon \mathrm{Pro-}\mathcal{C}(\pi) \to \mathcal{C}'(\pi)$，它给投射系 $Q = (Q_i)$ 对应 $\mathcal{C}'(\pi)$ 的对象 $X = \varprojlim_i Q_i$. 为定义反方向的函子，只需定义 $\mathcal{C}'(\pi)$ 到 $\mathcal{C}$-Ens 左正合函子范畴的反变函子，对每个 $X \in \mathcal{C}'(\pi)$ 取函子 $h(X)(E) = \operatorname{Hom}(X, E)$ (Hom 在 $\mathcal{C}'(\pi)$ 中取). 由定义立即看出函子 $h$ 和 $g$ 互为伴随，且 $hg$ 典范同构于 Pro-$\mathcal{C}(\pi)$ 的恒等函子. 剩下证明 (为建立 $g$ 和 $h$ 互为拟逆) $\mathcal{C}'(\pi)$ 的每个对象同构于形如 $g(Q)$ 的对象，其中 $Q \in \mathrm{Pro-}\mathcal{C}(\pi)$；换言之：每个以拓扑群 $\pi$ 为算子群、紧致且完全不连通的算子空间 $X$ 都同构于有限离散算子空间的投射极限. 由于 $X$ 是它的有限离散商 (作为无算子拓扑空间) 的投射极限，归结为证明在这些商的集合中存在 $\pi$-不变的共尾系. 为此只需证明：对这样的商 $X'$，$X'$ 在 $\pi$ 的作用下的变换集有限 (于是取这些变换的上确界，它是包含 $X'$ 的不变商). 或等价地，存在 $\pi$ 的开不变子群 $\pi'$，使该子群的元素不变 $X'$. 而 $X'$ 对应于 $X$ 的有限开集 $X_i$ 划分. 由 $\pi$ 的连续性和紧致性，存在单位元邻域 $V$，使得 $s \in V$ 蕴含对每个 $i$ 有 $s \cdot X_i \subset X_i$，故 $s$ 不变 $X'$. 而知道 $\pi$ 的开不变子群构成单位元的基本邻域组，证明完毕.
\end{proof}

注意还可以更简单地看到 Ind-$\mathcal{C}(\pi)$ 范畴典范等价于 $\pi$ 在其上连续作用的集合范畴. 我们在此用不到.

\begin{proposition}\label{prop:V.5.3}
设 $\mathcal{C}$ 是 Galois 范畴，$F$ 是 $\mathcal{C}$ 上的基本函子，$P = (P_i)$ 是按通常方式正规化的相伴 pro-对象. 设 $X \in \mathcal{C}$；为使 $X$ 连通，必要且充分条件是 $\pi$ 传递作用在 $E = F(X)$ 上.
\end{proposition}

\begin{proof}
归结到典型情形 $\mathcal{C} = \mathcal{C}(\pi)$，$F$ = 典范函子，在那里这是平凡的.
\end{proof}

\begin{corollary}\label{cor:V.5.4}
$X$ 的下列条件等价：(i) $X$ 连通且 $\neq \varnothing_{\mathcal{C}}$；(ii) 群 $\pi$ 在 $E = F(X)$ 上传递，且 $F(X) \neq \varnothing$；(iii) $X$ 同构于某个 $P_i$.
\end{corollary}

\begin{proof}
(i) 与 (iii) 的等价性也已经容易由第 4 节的 e) 得出.
\end{proof}

\begin{proposition}\label{prop:V.5.5}
设 $Q = (Q_i)_{i \in I}$ 是 $\mathcal{C}$ 的按通常方式正规化的 pro-对象，$G$ 是相应的 $\mathcal{C}$ 到 (Ens) 的函子 $G(X) = \operatorname{Hom}(Q, X)$. 下列条件等价：
\begin{enumerate}[(i)]
\item $G$ 与有限直和可交换；
\item $G$ 与两个对象的和可交换；
\item 诸 $Q_i$ 连通且 $\neq \varnothing$；
\item $Q$ 同构于 $\pi/H$，其中 $H$ 是 $\pi$ 的闭子群；
\item 函子 $G$ 同构于由 $\pi$ 的闭子群 $H$ 定义的函子 $E \mapsto E^{H}$ ($H$ 的不变元集).
\end{enumerate}
\end{proposition}

N.B. 在 (iv) 和 (v) 的陈述中，假设选定了基本函子，从而把 $\mathcal{C}$ 等同于范畴 $\mathcal{C}(\pi)$.

\paragraph*{证明.} 可以假设 $\mathcal{C} = \mathcal{C}(\pi)$. 蕴含 (i)$\Rightarrow$(ii) 平凡，(ii)$\Rightarrow$(iii) 像第 4 节的 d) 那样证明. 证明 (iii)$\Rightarrow$(iv). 事实上，知道 $\lim Q_i$ 作为非空有限集的投射极限非空. 设 $a$ 是 $\lim Q_i$ 的点，它定义算子空间同态
\[
\pi \longrightarrow Q,
\]
它满射，因为对每个 $i$，复合 $\pi \to Q \to Q_i$ 满射 (由 \ref{prop:V.5.3}，$\pi$ 在 $Q_i$ 中传递). 故若 $H$ 是 $a$ 的稳定化子群，得到同构 $\pi/H \xrightarrow{\sim} Q$. 蕴含 (iv)$\Rightarrow$(v) 和 (v)$\Rightarrow$(i) 又是平凡的.

\begin{proposition}\label{prop:V.5.6}
设 $\mathcal{C}$ 是 Galois 范畴，$P$ 是 $\mathcal{C}$ 的基本 pro-对象，$F$ 是相伴的基本函子. 设 $P' = (P'_i)_{i \in I}$ 是 $\mathcal{C}$ 的 pro-对象，化为标准形式，$F'$ 是相伴函子 $F'(X) = \operatorname{Hom}(P', X)$ ($\mathcal{C}$ 到 Ens). 下列条件等价：
\begin{enumerate}[(i)]
\item $P' \simeq P$，或 $F' \simeq F$；
\item $P'$ 基本，或 $F'$ 基本；
\item $F'$ 把两个对象的和变为和，且 $X \neq \varnothing_{\mathcal{C}}$ 蕴含 $F(X) \neq \varnothing$；
\item $\mathcal{C}$ 的连通且 $\neq \varnothing_{\mathcal{C}}$ 的对象恰是那些同构于某个 $P'_i$ 的对象.
\end{enumerate}
\end{proposition}

\begin{proof}
平凡地有 (i)$\Rightarrow$(iii) 和 (i)$\Rightarrow$(ii)，此外由 \ref{cor:V.5.4} (应用于 $P'$ 以代替 $P$) 有 (ii)$\Rightarrow$(iv). 再者，(iii) 或 (iv) 由 \ref{prop:V.5.5} 蕴含 $P'$ 形如 $\pi/H$，其中 $H$ 是 $\pi$ 的闭子群. 在情形 (iii)，对 $\pi$ 的每个开不变子群 $\pi'$ 存在 $\pi$-同态 $P' = \pi/H \to \pi/\pi'$，故 $H \subset \pi'$，从而 $H = (e)$，由此得 (i)，证毕.
\end{proof}

\begin{corollary}\label{cor:V.5.7}
设 $\mathcal{C}$ 是 Galois 范畴. 基本 pro-对象同构，基本函子同构.
\end{corollary}

\begin{proof}
换言之，基本函子范畴是连通\emph{群胚} (groupoid) $\Gamma$，可称为 Galois 范畴 $\mathcal{C}$ 的\emph{基本群胚} (fundamental groupoid)\index[term]{jibenqunpei@基本群胚}. 若 $\mathcal{C} = \mathcal{C}(\pi)$，则基本群胚的对象自同构群同构于 $\pi$，该同构在内自同构意义下完全确定. (N.B. 称所有态射都是同构的范畴为群胚，连通群胚即所有对象都同构的群胚.) $\mathcal{C}$ 的基本 pro-对象构成一个连通群胚，等价于基本群胚的对偶. 若 $F$，$F'$ 是两个基本函子，相伴的基本 pro-对象为 $P$，$P'$，则 $\mathrm{Hom}(F, F') = \mathrm{Isom}(F, F')$ 有时记为 $\pi_{F', F}$\index[notations]{$\pi_{F',F}$}\index[notations]{$\pi_{F',F}$}，扮演``从 $F$ 到 $F'$ 的\emph{道路类}集合''的角色. 特别地，$\pi_{F, F} = \pi_F$ 正是前一节构造的 $\mathcal{C}$ 在 $F$ 处的基本群. 至于与 $F$ 相伴的 pro-对象 $P$，它扮演 $\mathcal{C}$ 的终对象 $e_{\mathcal{C}}$ 在 $F$ 处的\emph{万有覆盖} (universal covering)\index[term]{wanyoufugai@万有覆盖} 的角色.
\end{proof}

给出 $\mathcal{C}$ (在等价意义下) 用其基本群胚 $\Gamma$ 的描述 (不经过选择其对象 $F$) 可能更方便. 对 $\mathcal{C}$ 的每个对象 $X$ 相伴基本群胚上的函子 $E_X$，定义为
\[
E_X(F) = F(X),
\]
取值于 Ens. (在拓扑中这样的函子称为群胚上的``\emph{局部系}'' (local system)\index[term]{jubuxi@局部系}.) $F(X) = E_X(F)$ 可称为 $X$ 在 $F$ 处的纤维 (fiber)，函子 $E_X$ 称为与 $X$ 相伴的纤维函子\index[term]{xianweihanzi@纤维函子}. 函子 $E_X$ 有下列性质：对每个 $F$，$E_X(F)$ 是拓扑群 $\pi_F = \operatorname{Aut}(F)$ 在其上连续作用的有限集. 对给定的从基本群胚到 Ens 的共变函子 $\xi$，前面的条件等价于对任意固定的 $F$ 的同样条件. 据此：

\begin{proposition}\label{prop:V.5.8}
函子 $X \mapsto E_X$ 是 $\mathcal{C}$ 与满足上面指出的条件的、$\mathcal{C}$ 的基本群胚 $\Gamma$ 到 Ens 的共变函子范畴之间的等价.
\end{proposition}

\begin{proof}
事实上，设 $F_0$ 是基本群胚的对象，$\pi_0 = \mathrm{Aut}(F_0)$，则函子 $\xi \mapsto \xi(F_0)$ 是 \ref{prop:V.5.8} 所考虑的第二个范畴与范畴 $\mathcal{C}(\pi_0)$ 之间的等价，立即看出. 另一方面，后者与函子 $X \mapsto E_X$ 的复合是自然等价 $\mathcal{C} \to \mathcal{C}(\pi_0)$. 由此得出函子 $X \mapsto E_X$ 本身是等价.
\end{proof}

\begin{corollary}\label{cor:V.5.9}
范畴 Pro-$\mathcal{C}$ 典范等价于基本群胚 $\Gamma$ 到拓扑空间范畴的共变函子 $\xi$ 的范畴，这些函子满足条件：对 $\Gamma$ 的每个对象 $F$，$\xi(F)$ 是以拓扑群 $\pi_F$ 为算子群的紧致完全不连通算子空间.
\end{corollary}

\begin{proof}
这里再次可以对 $\xi$ 验证这个条件，只需对一个 $F$ 验证. 证明与 \ref{prop:V.5.8} 相同.
\end{proof}

\begin{remark}\label{rem:V.5.10}
设 $(F_s)_{s \in S}$ 是基本群胚 $\Gamma$ 的对象族. 对 $s, s' \in S$ 令
\[
\operatorname{Hom}(s, s') = \operatorname{Hom}(F_s, F_{s'}),
\]
于是 $S$ 本身成为连通群胚，映射 $s \mapsto F_s$ 是 $S$ 到 $\Gamma$ 的全忠实函子，记为 $f$. 考虑 $\mathcal{C}$ 到函子范畴 $\mathrm{Hom}(S, \mathrm{Ens})$ 的函子 $X \mapsto E_X \circ f$，得到 \ref{prop:V.5.8} (和 \ref{cor:V.5.9}) 以 $\Gamma$ 换成 $S$ 的变体. 当 $S$ 约化为一点时，如此得到的陈述归结为定理 \ref{thm:V.4.1}，当 $S$ 是 $\Gamma$ 的对象集时正是 \ref{prop:V.5.8} 本身.
\end{remark}

我们将用 \ref{cor:V.5.9} 定义 $\mathcal{C}$ 的典范 pro-对象. 为此考虑 $\Gamma$ 到拓扑空间 (甚至拓扑群) 范畴的函子：
\[
f \colon F \longmapsto \operatorname{Aut}(F) = \pi_F.
\]
该函子满足 \ref{prop:V.5.8} 所考虑的条件，算子空间 $f(F)$ (在 $\pi_F$ 下) 正是 $\pi_F$ 本身，看作在自身内自同构作用下的算子空间. 故函子 $f$ 对应 $\mathcal{C}$ 的一个在唯一同构意义下确定的 pro-对象，它甚至是 $\mathcal{C}$ 的 pro-群，称为 $\mathcal{C}$ 的\emph{基本 pro-群} (fundamental pro-group)\index[term]{jiben pro-qun@基本 pro-群}，扮演\emph{基本群局部系}的角色. 因此它是 $\mathcal{C}$ 的 pro-群 $\Pi$，由下述条件定义：有对 $F$ 函子性的同构
\[
F(\Pi) \simeq \pi_F.
\]
若 $X$ 是 $\mathcal{C}$ 的任意 pro-对象，有典范态射
\[
\Pi \times X \longrightarrow X,
\]
它使 $X$ 成为 Pro-$\mathcal{C}$ 中以 $G$ 为左算子群的算子对象. 只需注意对变动的 $F$，有典范映射
\[
\Pi(F) \times X(F) \longrightarrow X(F),
\]
i.e.
\[
\operatorname{Aut}(F) \times E_X(F) \longrightarrow E_X(F), \quad \text{或} \quad \pi_F \times F(X) \longrightarrow F(X),
\]
它关于 $F$ 函子性. 它也关于 $X$ 函子性，故对 pro-对象的每个态射 $X \to Y$，相应的图交换.

\begin{remark}\label{rem:V.5.11}
务必不要把基本 pro-对象 $P$ (它不配备群结构，且连通) 与基本 pro-群 (它是 pro-群，一般不必连通) 混淆. 精确地说，$G$ 连通当且仅当 $\pi_F$ 通过内自同构作用在自身上传递，i.e. 当且仅当 $\pi$ 约化为单位元，或等价地 $\mathcal{C}$ 等价于有限集范畴. 另一个本质区别是 $G$ 在唯一同构意义下确定，而 $P$ 只在非唯一同构意义下确定.
\end{remark}

设 $E$ 是有限集，考虑群胚 $\Gamma$ 上取常值 $E$ 的函子：由 \ref{prop:V.5.8} 它定义 $\mathcal{C}$ 的对象，记为 $E_{\mathcal{C}}$，也可以解释为 $\mathcal{C}$ 的终对象 $e_{\mathcal{C}}$ 的 $E$ 份拷贝的和. 可以把 $E_{\mathcal{C}}$ 看作有限集范畴到范畴 $\mathcal{C}$ 的函子，该函子正合，故把有限群变为 $\mathcal{C}$ 的群，等等. 故若 $X$ 是有限群 $G$ 在其上右作用的 $\mathcal{C}$ 的对象，看到可以把 $X$ 看作具有右 $\mathcal{C}$-算子群 $G_{\mathcal{C}}$ 的 $\mathcal{C}$ 对象. 因此，按 $\mathcal{C}$-算子群对象一般术语的延伸，若 $X$ 在 $G_{\mathcal{C}}$ 下形式主齐性，i.e. 若由 $G_{\mathcal{C}}$ 在 $X$ 上的右作用得出的典范态射
\[
X \times G_{\mathcal{C}} \longrightarrow X \times X
\]
是同构，则称 $X$ 在 $G$ 下形式主齐性. 称 $X$ 在 $G$ 下主齐性\index[term]{zhuqixingduixiang@主齐性对象}，若它在 $G_{\mathcal{C}}$ 下主齐性，i.e. 若它形式主齐性，且此外商 $X/G = X/G_{\mathcal{C}}$ 是 $e_{\mathcal{C}}$. 若选定基本函子，从而给出 $\mathcal{C}$ 与某个范畴 $\mathcal{C}(\pi)$ 的等价，$X$ 对应一个 $\pi$ 在其上连续左作用的集合，即 $E = F(X)$. 让 $G$ 右作用在 $X$ 上，等价于让 $G$ 右作用在集合 $E$ 上，使 $G$ 的作用与 $\pi$ 的作用交换. 立即看到 $X$ 在 $G$ 下主齐性当且仅当集合 $E$ 是 $G$ 下的主齐性空间，i.e. 当且仅当 $G$ 在其上简单传递作用. (此外，$X$ 形式主齐性当且仅当 $E$ 主齐性或空.) 与 \ref{prop:V.5.3} 比较，看到若 $X$ 在 $G$ 下主齐性且连通，则给定同态 $G \to \operatorname{Aut}(X)$ 的对偶群是同构；且为使 $\mathcal{C}$ 的对象 $X$ 连通且在 $\mathrm{Aut}(X)$ 的对偶群下主齐性，必要且充分条件 (用第 4 节的记号) 是它同构于某个 Galois 的 $P_i$. 在典型情形 $\mathcal{C} = \mathcal{C}(\pi)$，这意味着 $X$ 同构于 $\pi$ 除以一个不变子群的商.

仍假设给定基本函子 $F$. 则给出在有限群 $G$ 下主齐性、带 $a \in F(X)$ 的 $X$，等价于给出从 $\pi$ 到群 $G$ 的同态. 事实上，这样的同态给出集合 $E = G$，让 $\pi$ 借助给定同态 $\pi \to G$ 和 $G$ 的左平移在其上左作用，让 $G$ 通过右平移右作用，$E$ 的标记点 $a$ 是 $G$ 的单位元. 由前面所述，由此以基本唯一的方式得到具有上述性质的每个三元组 $(X, G, a)$，因为带标记点的 $G$ 下主齐性集合等同于 $G$ 本身. 这样，有限群范畴到 Ens 的函子 $G \mapsto \operatorname{Hom}(\pi, G)$ 有直接的几何解释，该函子借助 $\pi$ pro-可表示，对它的考察因此给出与 $F$ 相伴的群 $\pi$ 的另一个构造.

\section{从一个 Galois 范畴到另一个的正合函子}

\begin{proposition}\label{prop:V.6.1}
设 $\mathcal{C}$，$\mathcal{C}'$ 是两个 Galois 范畴，$H \colon \mathcal{C} \to \mathcal{C}'$ 是共变函子，$F'$ 是 $\mathcal{C}'$ 上的基本函子，$F = F' \circ H$. 下列条件等价：
\begin{enumerate}[(i)]
\item $H$ 正合，i.e. 左正合且右正合；
\item $H$ 左正合，把有限和变为有限和，且把满态射变为满态射 (或：把 $\neq \varnothing_{\mathcal{C}}$ 的对象变为 $\neq \varnothing_{\mathcal{C}'}$ 的对象)；
\item $F$ 是 $\mathcal{C}$ 上的基本函子.
\end{enumerate}
\end{proposition}

\begin{proof}
蕴含 (i)$\Rightarrow$(ii) 是范畴论的一般事实. 此外 (ii) 的第一种形式蕴含第二种，这注意到：若 $X$ 是 $\mathcal{C}$ 的对象，则 $X \neq \varnothing_{\mathcal{C}}$ 当且仅当态射 $X \to e_{\mathcal{C}}$ 是满态射；注意 $F$ 假设左正合，把 $e_{\mathcal{C}}$ 变为 $e_{\mathcal{C}'}$. (ii) 的第二种形式蕴含 (iii)，因为 $F$ 左正合从而 pro-可表示，可以套用判则 \ref{prop:V.5.6} (iii). 最后 (iii) 蕴含 (i)，由 $F$ 正合且``保守'' (i.e. 满足第 4 节的公理 (G 6)) 的事实得出.
\end{proof}

设 $\Gamma$ 是 $\mathcal{C}$ 的基本群胚，$\Gamma'$ 是 $\mathcal{C}'$ 的基本群胚. 故若 $H$ 正合，则 $F' \mapsto F' \circ H$ 是从群胚 $\Gamma'$ 到群胚 $\Gamma$ 的函子，记为 ${}^{t}H$：
\[
{}^{t}H(F')(X) = F'(H(X)),
\]
它也可以用第 6 节引入的记号 $F(X) = E_X(F)$ 写成
\[
E_{H(X)}(F') = E_X({}^{t}H(F')).
\]
后一个公式 (考虑到 \ref{prop:V.5.8} 或 \ref{thm:V.4.1}) 表明正合函子 $H$ 在知道相应函子 ${}^{t}H$ 时 (在唯一同构意义下) 被确定. 固定 $F'$，令 $F = {}^{t}H(F')$，则 ${}^{t}H$ 定义同态 $\pi_{F'} = \operatorname{Aut}(F') \to \pi_F = \operatorname{Aut}(F)$：
\[
{}^{t}H \colon \pi_{F'} \longrightarrow \pi_F \quad (F = {}^{t}H(F') = F' \circ H).
\]
此外，上面的公式 (考虑到 \ref{prop:V.5.8}) 表明这个同态具有下列性质：对每个 $\pi_F$ 在其上连续作用的有限集 $E$，群 $\pi_{F'}$ 借助前面的同态 $\pi_{F'} \to \pi_F$ 也连续作用. 把它应用于 $\pi_F$ 除以开不变子群的商，看到前面的条件也意味着所考虑的同态连续. 反过来，给出 $\Gamma$ 的对象 $F$，$\Gamma'$ 的对象 $F'$ 和连续同态
\[
u \colon \pi_{F'} \longrightarrow \pi_F,
\]
它就给出从 $\mathcal{C}(\pi)$ 到 $\mathcal{C}(\pi')$ 的函子，显然正合，故由 \ref{thm:V.4.1} 它给出从 $\mathcal{C}$ 到 $\mathcal{C}'$ 的正合函子 $H$，使得 ${}^{t}H \colon \pi_{F'} \to \pi_F$ 恰为 $u$. 也可以不从群同态出发，而从函子
\[
U \colon \Gamma' \longrightarrow \Gamma
\]
出发，它满足：对每个 $F' \in \Gamma'$ (或一个 $F' \in \Gamma'$，都一样)，相应同态 $\pi_{F'} \to \pi_F$ 连续：这样的函子同构于形如 ${}^{t}H$ 的函子，其中 $H \colon \mathcal{C} \to \mathcal{C}'$ 是在唯一同构意义下确定的正合函子. 因此：

\begin{corollary}\label{cor:V.6.2}
为使 Galois 范畴的函子 $H \colon \mathcal{C} \to \mathcal{C}'$ 正合，必要且充分条件是存在等价 $\mathcal{C}(\pi) \to \mathcal{C}$ 和 $\mathcal{C}' \to \mathcal{C}(\pi')$，把函子 $H$ 变为与拓扑群同态 $\pi' \to \pi$ 相伴的函子 $\mathcal{C}(\pi) \to \mathcal{C}(\pi')$.
\end{corollary}

\begin{corollary}\label{cor:V.6.3}
设 $\mathcal{C}$，$\mathcal{C}'$ 是两个 Galois 范畴，$\Gamma$，$\Gamma'$ 是它们的基本群胚. 则从 $\mathcal{C}$ 到 $\mathcal{C}'$ 的正合函子范畴等价于具有下列性质的函子 $U \colon \Gamma' \to \Gamma$ 的范畴：对 $\Gamma'$ 中的每个 $F'$ (或 $\Gamma'$ 中的一个 $F'$，都一样)，令 $\mathcal{F} = \mathcal{U}(\mathcal{F}')$，由 $U$ 定义的同态
\[
\pi_{F'} = \operatorname{Aut}(F') \longrightarrow \pi_F = \operatorname{Aut}(F)
\]
连续.
\end{corollary}

\begin{proof}
考虑 $\mathcal{C}$ 的基本 pro-群 $\Pi$，正合函子 $H$ 把它变为 $\mathcal{C}'$ 的 pro-群 $H(\Pi)$，定义同态
\[
\Pi' \longrightarrow H(\Pi)
\]
(其中 $\Pi'$ 是 $\mathcal{C}'$ 的基本 pro-群)，其条件为：对 $\Gamma'$ 的每个对象 $F'$，相应的同态
\[
F'(\Pi') = \pi_{F'} \longrightarrow F'(H(\Pi)) = \pi_F \quad \left(\text{其中 } F = F' \circ H = {}^{t}H(F')\right)
\]
是自然同态
\[
\operatorname{Aut}(F') \longrightarrow \operatorname{Aut}(F' \circ H).
\]
由于后者关于 $F'$ 函子性，由 \ref{prop:V.5.8} 它确实定义 $\mathcal{C}'$ 的 pro-对象、事实上 pro-群的同态. 称该同态为与函子 $H$ 相伴的同态.
\end{proof}

现在设 $H'$ 是从 Galois 范畴 $\mathcal{C}'$ 到 Galois 范畴 $\mathcal{C}''$ 的第二个正合函子. 平凡地有
\[
{}^{t}(H'H) = {}^{t}H\,{}^{t}H'
\]
(N.B. 这里是函子的恒等，而不只是典范同构). 基本 pro-群的相伴同态有类似的传递性.

现在用相应同态
\[
u \colon \pi_{F'} \longrightarrow \pi_F \quad (\text{其中 } F' = F' \circ H)
\]
解释正合函子 $H$ 的性质.

引入 (配备基本函子 $F$ 的) Galois 范畴 $\mathcal{C}$ 的\emph{带点对象} (pointed object) 概念是方便的\index[term]{daidianduixiang@带点对象}：它按定义是配备 $F(X)$ 中元素 $a$ 的 $\mathcal{C}$ 的对象 $X$. 它因此解释为 $\pi_F$ 在其上连续左作用的有限集，配备一点 $a$. 故由 \ref{prop:V.5.3}，$\mathcal{C}$ 的连通带点对象等同于 $\pi_F$ 的开子群. 若 $U$ 和 $V$ 是两个这样的子群，对应于 $\mathcal{C}$ 的连通带点对象 $X$，$Y$，则存在从 $X$ 到 $Y$ 的带点同态当且仅当 $U \subset V$，且此时该同态唯一. 当然，函子 $H$ 把带点对象变为带点对象 (因为 $F = F' \circ H$). 另一方面注意，像 $\pi_F$ 这样的群的闭子群是包含它的开子群的交；因此，若 $M$，$N$ 是两个闭子群，则 $M \subset N$ 当且仅当每个包含 $N$ 的开子群也包含 $M$. 借助这些评注，容易证明下列结果：

\begin{proposition}\label{prop:V.6.4}
设 $X$ 是 $\mathcal{C}$ 的连通带点对象，对应于 $\pi_F$ 的开子群 $U$. 为使 $U$ 包含 $u(\pi_{F'})$ (相应地：$u(\pi_{F'})$ 生成的闭不变子群)，必要且充分条件是 $H(X)$ 有带点截面 (相应地：完全分裂).
\end{proposition}

称对象 $X$ (在终对象之上的) 截面 (section) 为 Galois 范畴 $\mathcal{C}$ 的对象 $X$ 的从终对象 $e_{\mathcal{C}}$ 出发的态射，它等价于给出 $\pi_F$ 不变的 $F(X)$ 的元素 $a$；若 $X$ 带点，当截面与 $X$ 和 $e_{\mathcal{C}}$ 上的带点结构相容，i.e. 当 $a$ 恰是 $F(X)$ 的标记对象时，称有带点截面 (pointed section). 这样的截面因此唯一，且存在当且仅当 $F(X)$ 的标记对象被 $\pi_F$ 不变. 最后，若 Galois 范畴的对象同构于若干终对象的和，i.e. 若 $\pi_F$ 在 $F(X)$ 中平凡作用 —— 当 $X$ 带点时，这显然是比存在带点截面更强的条件 —— 则称该对象\emph{完全分裂} (completely decomposed)\index[term]{wanquanfenlie@完全分裂}. 命题 \ref{prop:V.6.4} 由定义和前面的评注平凡地得出.

\begin{corollary}\label{cor:V.6.5}
为使 $u$ 平凡，必要且充分条件是：对 $\mathcal{C}$ 的每个对象 $X$，$H(X)$ 完全分裂.
\end{corollary}

\begin{proposition}\label{prop:V.6.6}
设 $X'$ 是 $\mathcal{C}'$ 的连通带点对象，对应于 $\pi_{F'}$ 的开子群 $U'$. 为使 $U'$ 包含 Ker $u$，必要且充分条件是存在 $\mathcal{C}$ 的连通带点对象 $X$ 和从 $H(X)$ 的连通带点分量 $X'_0$ 到 $X'$ 的带点同态 (故 $X$ 作为带点对象同构于 $H(X)$ 的原像的带点中性连通分量的商). 若 $u$ 满射，前面的条件也等价于：$X'$ 同构于某个 $H(X)$，其中 $X$ 是 $\mathcal{C}$ 的带点对象.
\end{proposition}

(Galois 范畴 $\mathcal{C}$ 的带点对象 $X$ 的中性连通分量 (neutral connected component)，是 $X$ 的惟一的连通带点子对象；由 \ref{prop:V.5.3}，它对应于 $\pi_F$ 在 $F(X)$ 的标记点上的轨道.) 由于``$U'$ 包含 Ker $u$''不依赖于 $X'$ 所选择的带点结构 (因为另一个带点结构就是把 $U$ 换成共轭子群 $\dot{\mathrm{a}} U$)，看到：

\begin{corollary}\label{cor:V.6.7}
为使 $U'$ 包含 Ker $u$，必要且充分条件是存在 $\mathcal{C}$ 的对象 $X$ (可以假设连通) 和从 $H(X)$ 的一个连通分量到 $X'$ 的态射. 若 $u$ 满射，这也意味着 $X'$ 同构于形如 $H(X)$ 的对象.
\end{corollary}

\begin{corollary}\label{cor:V.6.8}
为使 $u$ 单射，必要且充分条件是：对 $\mathcal{C}'$ 的每个对象 $X'$，存在 $\mathcal{C}$ 的对象 $X$ 和从 $H(X)$ 的一个连通分量到 $X'$ 的同态.
\end{corollary}

\begin{proposition}\label{prop:V.6.9}
下列条件等价：
\begin{enumerate}[(i)]
\item 同态 $u \colon \pi_{F'} \to \pi_F$ 满射；
\item 对 $\mathcal{C}$ 的每个连通对象 $X$，$H(X)$ 连通；
\item 函子 $H$ 全忠实.
\end{enumerate}
\end{proposition}

\begin{proof}
最后一个事实意味着对 $\mathcal{C}$ 的两个对象 $X$，$Y$，自然映射
\[
\operatorname{Hom}(X, Y) \longrightarrow \operatorname{Hom}(H(X), H(Y))
\]
双射.
\end{proof}

\begin{corollary}\label{cor:V.6.10}
为使 $u$ 是同构，必要且充分条件是 $H$ 是范畴等价，或等价地下列两个条件成立：
\begin{enumerate}[a)]
\item 对 $\mathcal{C}$ 的每个连通对象 $X$，$H(X)$ 连通；
\item $\mathcal{C}'$ 的每个对象同构于形如 $H(X)$ 的对象.
\end{enumerate}
\end{corollary}

\begin{proposition}\label{prop:V.6.11}
设 $H \colon \mathcal{C} \to \mathcal{C}'$ 和 $H' \colon \mathcal{C}' \to \mathcal{C}''$ 是 Galois 范畴之间的正合函子，$F''$ 是 $\mathcal{C}''$ 上的基本函子，令 $F' = F''H'$ 和 $F = F'H$，考虑相伴同态
\[
u' \colon \pi_{F''} \longrightarrow \pi_{F'} \qquad u \colon \pi_{F'} \longrightarrow \pi_F.
\]
为使 Ker $u \subset$ Im $u'$ (i.e. 为使 $uu'$ 是平凡同态)，必要且充分条件是：对 $\mathcal{C}$ 的每个对象 $X$，$H'(H(X))$ 完全分裂. 为使 Ker $u \supset$ Im $u'$，必要且充分条件是：对 $\mathcal{C}'$ 的每个使得 $H'(X')$ 有带点截面的连通带点对象 $X'$，存在 $\mathcal{C}$ 的对象 $X$ 和从 $H(X)$ 的一个连通分量到 $X'$ 的同态.
\end{proposition}

\begin{proof}
第一个断言由 \ref{prop:V.6.4} 的最后一个断言得出. 第二个断言由 \ref{prop:V.6.4} 和 \ref{prop:V.6.6} 的结合得出.
\end{proof}

\begin{remark}\label{rem:V.6.12}
一般地，在 \ref{cor:V.6.8} 的条件下，$X'$ 同构于形如 $H(X)$ 的对象并不成立. 可以证明，为使 $\mathcal{C}'$ 的每个连通 (从而每个) 对象同构于形如 $H(X)$ 的对象，必要且充分条件是 $u$ 是 $\pi_{F'}$ 到 $\pi_F$ 的一个直因子子群上的同构. 在实践中，用 $\mathcal{C}'$ 到 $\mathcal{C}$ 的适当正合函子直接构造 $u$ 的右逆 $\pi_F \to \pi_{F'}$.
\end{remark}

\begin{proposition}\label{prop:V.6.13}
设 $\mathcal{C}$ 是配备基本函子 $F$ 的 Galois 范畴，$S$ 是 $\mathcal{C}$ 的连通对象，$\mathcal{C}'$ 是 $\mathcal{C}$ 在 $S$ 之上的对象范畴. 则 $\mathcal{C}'$ 是 Galois 范畴，且从 $\mathcal{C}$ 到 $\mathcal{C}'$ 的函子 $X \mapsto H(X) = X \times S$ 正合. 设 $a \in F(S)$，设 $F'$ 是从 $\mathcal{C}'$ 到有限集范畴的函子，定义为
\[
F'(X') = a \text{ 在 } F(X') \longrightarrow F(S) \text{ 下的原像}.
\]
则有同构 $F \cong F' \circ H$，且相应同态
\[
u \colon \pi_{F'} \longrightarrow \pi_F
\]
是 $\pi_{F'}$ 到 $\pi_F$ 的稳定化子 $F(X)$ 的标记元素 $a$ 的开子群 $U$ 上的同构.
\end{proposition}

\begin{proof}
证明留给读者.
\end{proof}

\section{预概形的情形}

设 $S$ 是局部诺特连通预概形，设
\[
a \colon \operatorname{Spec}(\Omega) \longrightarrow S
\]
是 $S$ 的几何点 (geometric point)，取值于代数闭域 $\Omega$. 令
\[
\mathcal{C} = S \text{ 的平展覆盖的范畴}
\]
对 $\mathcal{C}$ 的对象 $X$，i.e. $S$ 的平展覆盖 $X$，令
\[
F(X) = X \text{ 在 } a \text{ 之上的几何点的集合}.
\]
于是 $F$ 成为 $\mathcal{C}$ 上取值于有限集范畴的函子. 性质 (G 1) 至 (G 6) 满足：对 (G 1)，它包含在 \ref{prop:I.4.6} 的平凡推论中；(G 2) 由 \ref{cor:V.3.4} 得出；(G 3) 由 \ref{prop:V.3.5} 得出；(G 4) 由定义平凡；(G 5) 由 \ref{prop:V.3.5} 和第 2 节开头得出；最后 (G 6) 在 \ref{prop:V.3.7} 中证明. 故可以应用第 4、5、6 节的结果. 这特别使我们能够定义表示 $F$ 的 $\mathcal{C}$ 的 pro-对象 $P$，称为 $S$ 在点 $a$ 的万有覆盖 (universal covering)\index[term]{wanyoufugai (yugaixing)@万有覆盖 (预概形)}，以及拓扑群 $\pi = \operatorname{Aut}(F) = \operatorname{Aut}^{0}P$，称为 $S$ 在 $a$ 处的基本群 (fundamental group)，记为 $\pi_1(S, a)$\index[notations]{$\pi_1(S, a)$}\index[notations]{$\pi_1(S, a)$}. 函子 $F$ 于是定义 $\mathcal{C}$ 与 $\pi = \pi_1(S, a)$ 在其上连续作用的有限集范畴的等价. 这个等价因此允许把覆盖上的通常投射极限和有限归纳极限运算 (积、纤维积、和、取商等) 解释为 $\mathcal{C}(\pi)$ 中的类似运算，i.e. 解释为 $\pi$ 在其上作用的有限集上的显然运算. 还注意，由于平展覆盖的拓扑连通分量也是平展覆盖，$\mathcal{C}$ 的对象 $X$ 在 $\mathcal{C}$ 中连通当且仅当 $X$ 拓扑连通；由 \ref{prop:V.5.3} 这意味着 $\pi_1$ 在 $F(X)$ 中传递作用. 注意，为使 $\mathcal{C}$ 的对象 $X$ 在 $S$ 上忠实平坦且拟紧 (它已在 $S$ 上平坦且拟紧)，必要且充分条件是 $X \to S$ 满射，i.e. 是 $\mathcal{C}$ 中的满态射，或 $X \neq \varnothing$. 于是由判则 \ref{prop:V.2.6} (iii) 得出结论：$X$ 是 $S$ 的以 $G$ 为 伽罗瓦群的主覆盖，当且仅当它是范畴 $\mathcal{C}$ 中 $G$ 下的主齐性空间 (如第 5 节所定义的).

若 $a'$ 是 $S$ 的另一个几何点 (对应于代数闭域 $\Omega'$，它可以不同于 $\Omega$，甚至可以有不同的特征)，它定义 $\mathcal{C}$ 到有限集范畴的纤维函子 $F' = F_{a'}$，它仍正合，故同构于 $F = F_a$. 因此，对变动的 $a$，基本群 $\pi_1(S; a)$ 彼此同构. 若用 $\pi_1(S; a, a')$ 记相伴纤维函子的同构\index[notations]{$\pi_1(S; a, a')$} (或等价地，同态) $F_a \to F_{a'}$ 的集合，则得到一个群胚，其对象集是 $S$ 的几何点集，基本群是所述群胚的对象的自同构群. 集合 $\pi_1(S; a', a)$ 可称为从 $a$ 到 $a'$ 的道路类集合\index[term]{daolulei@道路类}. 这些类因此以显然的方式复合. 最后，可以定义 $\mathcal{C}$ 的 pro-群 $\Pi_1^{S}$\index[notations]{$\Pi_1^S$}，可称为 $S$ 的基本 pro-群或 $S$ 上的基本群局部系\index[term]{jibenqunjubuxi@基本群局部系}，由下列条件在唯一同构意义下定义：有对 $S$ 的几何点 $a$ 函子性的同构
\[
F_a(\Pi_1^{S}) = \pi_1(S; a)
\]
(参见评注 \ref{rem:V.5.10}). 特别地，若 $s$ 是 $S$ 的普通点，则 $G$ 在 $s$ 处的纤维是 $k(s)$ 上的 pro-群，是 $k(s)$ 上有限平展群的投射极限；可以把这个 pro-群称为 $S$ 在普通点 $s$ 处的基本群，记为 $\pi_1(S, s)$. 按定义，$S$ 在 $k(s)$ 的代数闭扩张 $\Omega$ 中取值的这些点，正是 $\pi_1(S; a)$ 的元素，其中 $a$ 是由该扩张定义的 $S$ 的几何点. 特别地，(取 $S$ 为域的谱) 每个域 $k$ 典范地且函子性地相伴一个 $k$ 上的 pro-群，可记为 $\pi_1(k)$，它是 $k$ 上有限平展群的投射极限，其在 $k$ 的代数闭扩张 $\Omega$ 中的点等同于 $\overline{\{k\}}/k$ 的拓扑 伽罗瓦群 (其中 $\overline{k}$ 是 $k$ 在 $\Omega$ 中的 Galois 闭包) 的元素 (参见 \ref{prop:V.8.1}). 这个群 $\pi_1(k)$ 似乎尚未引起代数学家的注意.

现在设
\[
f \colon S' \longrightarrow S
\]
是从一个连通局部诺特预概形到另一个的态射，$a'$ 是 $S'$ 的几何点，$a = f(a')$ 是它在 $S$ 中的直接像. 则``原像''函子诱导从 $S$ 的平展覆盖范畴 $\mathcal{C}(S)$ 到 $S'$ 的平展覆盖范畴 $\mathcal{C}(S')$ 的函子：
\[
f^{\bullet} \colon \mathcal{C}(S) \longrightarrow \mathcal{C}(S').
\]
此外有函子同构
\[
F_a \cong F_{a'} \circ f^{\bullet},
\]
从而 $f^{\bullet}$ 是正合函子，可以对其应用第 6 节的结果. 特别有典范同态
\[
u = \pi_1(f; a')\colon \pi_1(S', a') \longrightarrow \pi_1(S; a) \quad (a' = f(a)),
\]
它可以重构原像函子，即作为算子群限制运算. 函子 $f^{\bullet}$ 的性质因此用相伴群同态的性质简单表达，如第 6 节所明确的. 若特别地 $S'$ 是 $S$ 的平展覆盖，则同态 $u$ 是 $\pi_1(S', a')$ 到定义 $S$ 的连通带点平展覆盖 $S'$ 的 $\pi_1(S, a)$ 的开子群 (i.e. $a' \in F_a(S')$ 在 $\pi_1(S; A)$ 中的稳定化子 $U$) 上的同构.

若想对变动的几何点 $a'$ 解释同态 $\pi_1(f; a')$，则应当 (按照第 6 节所说的) 考虑 pro-群同态
\[
\Pi_1(f) \colon \Pi_1^{S'} \longrightarrow f^{\bullet}(\Pi_1^{S}),
\]
并取几何纤维的相应同态.

\section{正规基预概形的情形}

\begin{proposition}\label{prop:V.8.1}
设 $S$ 是域 $k$ 的谱，$\Omega$ 是 $k$ 的代数闭扩张，定义 $S$ 在 $\Omega$ 中取值的几何点 $a$. 设 $\overline{k}$ 是 $k$ 在 $\Omega$ 中的可分闭包. 则存在 $\pi_1(S, a)$ 到 $\overline{k}/k$ 的拓扑 伽罗瓦群的典范同构.
\end{proposition}

\begin{proof}
设 $k'$ 是 $k$ 在 $\Omega$ 中的代数闭包，它对应 $S$ 在 $k'$ 中取值的几何点 $b$. 函子的自然同态 $F_b \to F_a$ 显然是同构，因为 $k$ 的有限可分扩张的 $k$-同态在 $\Omega$ 中取值必然在 $\overline{k}$ 中，更不用说在 $k'$ 中. 另一方面，$k$ 的 $k'/k$ 的自同构群 $\pi'$ 以显然的方式作用在 $F_b$ 上，由此得同态 $\pi' \to \operatorname{Aut}(F_b) \stackrel{\sim}{\longrightarrow} \operatorname{Aut}(F_a) = \pi_1(S; a)$. 另一方面，众所周知从 $\pi'$ 到 $\overline{k}/k$ 的自同构群 $\pi$ 的自然同态是同构. 由此得到典范同态 $\pi \to \pi_1(S; a)$，剩下证明它是同构. 事实上，该同态单射，因为核的元素是 $\overline{k}/k$ 的自同构，它在每个有限可分子扩张上诱导恒等，故平凡. 该同态满射，因为若 $X$ 是 $S$ 的连通平展覆盖，从而由 $k$ 的有限可分扩张 $L/k$ 定义，则 $\pi$ 在 $L$ 到 $k'$ 的 $k$-同态集合上传递，众所周知.
\end{proof}

\begin{proposition}\label{prop:V.8.2}
设 $S$ 是连通局部诺特正规预概形，$K = k(s)$ 是它的函数域 = 其在生成点 $s$ 处的剩余域，$\Omega$ 是 $K$ 的代数闭扩张，定义 $S' = \mathrm{Spec}(K)$ 的几何点 $a'$ 和 $S$ 的几何点 $a$. 则同态 $\pi_1(S'; a') \longrightarrow \pi_1(S; a)$ 满射. 当把第一个群等同于 $K$ 在 $\Omega$ 中的可分闭包 $\overline{K}$ 的 伽罗瓦群时 (参见 \ref{prop:V.8.1})，前面同态的核在 Galois 理论下对应于 $\overline{K}/K$ 的在 $S$ 上非分歧的 $K$ 的有限扩张组成的子扩张.
\end{proposition}

\begin{proof}
第一个断言意味着 $S$ 的连通平展覆盖 $X$ 在 $S'$ 上的原像连通，i.e. $X$ 整，这正是 (\ref{prop:I.10.1}). 前面同态的核解释为 $\overline{K}/K$ 的在集合 $F_a(X)$ 上诱导恒等的自同构 (可以假设 $S$ 的平展覆盖 $X$ 连通). 但这意味着该自同构在 $\overline{K}/K$ 的在 $S$ 上非分歧的有限子扩张上诱导恒等，这证明最后一个断言.
\end{proof}

\begin{remarkstar}
借助这种把正规预概形 $S$ 的基本群解释为通常 Galois 理论的方式，这个情形的定义很久以前就已知晓.
\end{remarkstar}

\section{非连通预概形的情形：多 Galois 范畴}

设 $S$ 是局部诺特预概形，$(S_i)_{i \in I}$ 是它的连通分量. 则 $S$ 的平展覆盖范畴 $\mathcal{C}(S)$ 等价于诸 $\mathcal{C}(S_i)$ 的乘积范畴，因此可以用诸 $S_i$ 的基本群 (在每个 $S_i$ 中选定一个几何点后) 描述. 在平展态射下降理论的应用中，有时为每个 $S_i$ 选择几何点很不方便. 更方便的是借助 \ref{prop:V.5.8} 的自然推广，把 $\mathcal{C}(S)$ 解释为 $S$ 的几何点群胚上的函子范畴，该群胚看作 $S$ 的诸连通分量相应群胚的和；所考虑的函子是取值于有限集范畴的函子，满足与 \ref{prop:V.5.8} 中所援引的类似的连续性性质. 在实践中，会有 $S$ 的几何点族 $(a_t)_{t \in E}$，使得 $S$ 的每个连通分量 $S_i$ 都含有一个这样的点，于是像 \ref{rem:V.5.10} 中那样，可以用底集为 $E$ 的类似群胚代替 $S$ 的所有几何点的群胚. 当然，这些考虑应当纳入等价于 $\mathcal{C}(\pi)$ 型范畴的乘积范畴的范畴的一般定义中，可称之为\emph{多 Galois 范畴} (multigaloisian categories)\index[term]{duo Galois fanchou@多 Galois 范畴}. 我们把细节留给读者.

